\section[第六章\q 系统之路]{第六章\\ 系统之路}

\subsection{贪婪与恐惧}

{\kaishu
在市场贪婪的时候要恐惧，在市场恐惧的时候要贪婪。

\hfill ——巴菲特

贪婪与恐惧来自于人类正常的本能，而不应该成为一切错误的替罪羔羊。

\hfill ——金融帝国
}

几乎所有的人都认为贪婪与恐惧是交易者的天敌。这种说法看似非常有道理，毕竟我们似乎都经历过由于对大利润的贪婪而错失了最佳离场时机，也经历过由于对亏损的恐惧而使我们在最后一跌中丢掉了持有的筹码。我希望读者朋友们能够深刻地思考一下贪婪与恐惧的本质，而不要将其当做所有“失误”的替罪羔羊。

贪婪与恐惧形成于人类祖先丛林生活中产生的原始本能。恐惧来源于“或战或逃反应”，如果我们祖先对于所有的危险都勇往直前的话，那么人类恐怕就无法繁衍至今天。贪婪来源于，生物在顺境中追求更多当时并不需要的生存资源而备逆境中使用，否则很多物种无法熬过寒冷的严冬。在我看来，贪婪与恐惧是我们繁衍至今必不可少的两种心理反应，并且他们早已深入我们的本能与潜意识，从而不可能被我们消除。对于我们这些远离丛林的现代人而言，贪婪与恐惧会表现得更加复杂。恐惧也许并非完全来自于损失，在大牛市中害怕踏空的恐惧往往更能表现得淋漓尽致；而贪婪也会发生在大熊市中，表现为一种希望以更低价格买入的贪婪。从某种意义上来讲，在特定的行情阶段中，也可能表现为一部分人的贪婪，与另一部分人的恐惧。

如果我们继续深入地研究，还能够发现很多时候恐惧反应表现得并不客观。有时候恐惧完全来自于“对过去痛苦的记忆”，这就是“一朝被蛇咬、十年怕井绳”的道理，而现实并没有出现任何危机的迹象。更有意思的是，根据“焦虑原理”，人们对不确定性的恐惧，往往比确定无疑的损失更为可怕。可能亏损一定数量的金钱，要比确实亏损一定数量的金钱更容易使交易者恐惧。由此可见，恐惧通常并非对现实的反应，而是一种对未来的担忧。

我在市场上，经常听到某某股民指着套牢惨重的股票说：“怪我太贪了，明明给了我百分之几的利润我都不跑，结果损失这么多。”言外之意就是，损失完全来自于他没有兑现百分之几的利润。在他说这句话的时候，是否真能在以后放弃超过百分之几的利润吗？亏损不是来自于交易策略的问题，不是来自没有止损的问题，而只是来自于贪婪吗？如果交易者总是拿贪婪与恐惧当做错误的替罪羔羊，而不进行任何反思的话，那么这个交易者注定一事无成！

不贪婪、不恐惧，我们真的能够做到吗？“不恐惧”就会让我们对所有危机勇往直前，是否就注定我们早晚会被市场扫地出门呢？错误难道真的不在于自身“或战或逃反应”的尺度并不客观，而只是恐惧的问题吗？我们来到股市都是抱着获利的动机，并且使用我们当时并不需要的初始资金。如果我们真的做到绝对不贪婪，那么我们就早已成佛，还来股市干什么呢？要知道交易利润并非我们生存的必要条件，那么我们进入股市的动机就完全来自于贪婪。反过来想想，如果我们不在顺境中积攒足够的利润，那么我们又如何应对交易中必然会出现的逆境呢？

按照“鸵鸟现象”的说法，对于我们无法改变的现实，最好的办法就是接受它、合理控制它。无论是本能还是动机，都注定了我们永远无法消除贪婪与恐惧的心理现象。那么为什么我们不能接受它、并且客观地量化它（将其固定在一定的尺度上）呢？

巴菲特说：“在市场贪婪的时候要恐惧，在市场恐惧的时候要贪婪。”由此可见，并非贪婪与恐惧的错误，而是人们在最不恰当的时候出现这两种心理现象。如果说，某交易者当亏损超过50％以上时才会出现恐惧，那么他的恐惧尺度绝对有问题！既然我们永远无法消除恐惧、甚至恐惧根本就不应该被消除，那么就客观地量化恐惧吧。事实上，止损行为就是一种恐惧的表现，他在潜在的风险面前选择了逃跑。孙子说过善败者不亡，那么勇往直前的匹夫灭亡只是早晚的事情。

按照“树形原理”的说法，同样的原因可能出现不同的结果。或者说，没有出现亏损的交易不等于不可能出现亏损。我觉得交易者应该放弃“不劳而获”的幻觉，树立一种成本观念。如果你从来都不会亏损，那么捕捉多小的利润都是可以接受的。交易获利的关键就是，在顺境中赚取足够逆境中亏损的金钱。即使从事后来看，你的某笔交易是获利的，但站在事先的角度来说，你这笔交易也是可能出现亏损的。在我看来，你顺、就是顺！根本没有必要限制自己的好运气。我们根本没有必要在收获的秋季放弃准备过冬的食物，而又在冰天雪地的严冬出外捕猎。我的意思是，贪婪并不是罪过，更不应该成为亏损的替罪羔羊！

这样一来，对恐惧的控制就自然形成一种趋势跟踪系统，与此同时我们还要将贪婪发挥到极致。如果承认任何交易方法都是基于概率上的产物，并且任何方法都是可能产生亏损的。那么我们的当务之急就是，不要继续在混乱中寻找并不存在的完美与确定，而要给自己的交易方法找一个固定的“家”。我之所以在交易系统章节的开头说这些，是希望大家能够接受一些我们无论如何都无法回避的成本，从而将自己的交易方法固定下来。这时一个可量化的交易系统就出现了。

\subsection{趋势是你的朋友}

{\kaishu
当你指责“趋势”没有尽到朋友的义务时，请先问问自己，是否真的拿“趋势”当朋友了？

\hfill ——金融帝国
}

“趋势是你的朋友”，这是每一位交易者都耳熟能详的名言。我在一场大型股指期货报告会上，听到某资深分析师引用此名言。结果是迷惑了我几分钟后，让我真的感到哭笑不得。我想大家都能理解“趋势”的含义，这里我主要和大家谈谈“朋友”的概念吧。

第一，对朋友有所求并不可耻。无论是我们小时候追求快乐的玩伴，还是长大以后心灵交流与寻求安慰的好朋友，或者是生意场上追求双赢的合作伙伴，这些都能算得上朋友。事实上，无论是为了追求快乐，还是追求心理安慰，或者是追求物质利益。朋友之间都是“有所图”的，也许这种“所图”只不过是一种开心、轻松的感觉罢了。

第二，不肯对朋友付出才是可耻的。朋友通常都是相互的，他是你的好朋友，那么你也是他的好朋友。如果朋友之间都是只想索取，而不肯付出的话，那么这种关系绝对不可能长久维系。希望天天占便宜而吃不得一点亏的朋友，你觉得可交吗？

第三，朋友之间最需要的是尊重。朋友之间应该是平等的，而不存在谁高谁一头的问题。如果你总是想证明自己比朋友更聪明，总是不断地表现自己比朋友更正确，甚至总想不断嘲笑你的朋友、把他当做傻子，那么你的朋友绝对会对你厌恶之极，他会如何对你也就可想而知了。

第四，朋友与路人的区别在于朋友更能看长。如果你与陌生的路人处事，可能就是一种同一时点上的纯利益交换。但真正的好朋友，并不用为每一次的得失而斤斤计较，而是用一种长期双赢的方式来交往。“现用人、现为人”的方式，绝对不是什么交友之道。

好了，结论出来了：你对趋势有所求并不可耻，但不能为趋势有丝毫的付出才是可耻的。如果你真想和趋势做朋友的话，那么你就应该尊重它。而不能总想着自作聪明地超越它，更不能对它呼来唤去。如果你总是对朋友呼来唤去，那么可想而知朋友会如何对待你。如果你非要和你的趋势朋友斤斤计较一城一池的得失，那么你的趋势朋友就会离你而去，因为它看不惯你“现用人、现为人”的处事风格。

那么对于更多的交易者而言，又究竟是想和趋势维持一种什么样的关系呢？这里我们可以做一个不太恰当的比喻：如果一位男孩同一位女孩经常接触的话，那么你会觉得他们是一种什么关系呢？也许是恋人、也许是朋友，但答案绝对不只是简单的二元选择。如果一位女孩既接受不了成为恋人所承受的弊端，又不想放弃被人追求所带来的心理满足，那么她真正想要的就既不是恋人关系、也不是朋友关系，她真正想要的只不过是一位“追求者”罢了。当然，这位女孩只能美其名曰地说：我只是想和你做朋友。问题在于，如果男孩真的能够坦然地拿她当朋友，那么恐怕她会发觉朋友关系也并非是她所想要的。事实上，更多的交易者同样也是既不想承担趋势跟踪的代价、也不想放弃趋势跟踪所能够带来的利益。这些自负的交易者，恐怕只是希望把趋势当做自己的“追求者”，同时还自认为是在同趋势做朋友。问题是，这种畸形的朋友关系是不可能长久维系的。发现没有？人在处理任何有关得失的问题时，都会表现出惊人相似的倾向，而很少有人能够真正懂得取舍的道理。

和趋势做朋友，你真的想好了吗？没有人强迫你和趋势做朋友，但如果你们真的做了朋友，那么在你希望享受朋友权利的时候，问问自己，是否尽到了朋友的义务呢？其实趋势这位朋友真的很单纯、真的很慷慨、真的很仗义。它会为真正的朋友付出大量的钞票，但它又会经常让朋友付出几次小小的成本，从而看看这位朋友是否对它真心。哈哈，你不想当傻子，趋势更不想当傻子。趋势可不想成为付出了大量的钞票，还被“朋友”嘲笑为傻子的“冤大头”。

\subsection{小偷的下场}

{\kaishu
我们从来都不曾真正地得到什么，我们得到的只不过是靠失去另一些东西而换回来的。

\hfill ——金融帝国
}

小偷是一种令人深恶痛绝的“职业”，他们总是非法地占有别人的劳动成果。从某种意义上来讲，小偷是那种希望不用付出就能够得到回报的人。抛开道德是非不谈，偷一次就洗手不干的贼还是有可能真正占到便宜的。但想靠偷东西为生的贼，被“剁手”只是早晚的事情！由此看来，做贼不是一种可以当做职业的行当。

律师和医生都要付出十多年的时间来学习专业知识，此后才能获得并不算庞大的收入。而一位交易者上来就想在市场上大赚钞票，是否就显得很可笑呢？交易是一种专业性很强的职业，甚至它的成功率都是极低的。无论你是否愿意，你都会付出并不亚于律师和医生学习专业知识时所花费的时间与金钱。他们不同之处在于，律师和医生出师后只能赚到有数的收入，而一个交易者则可能通过复利效应赚到天文数字般的财富。另外，律师与医生是否能够学有所成，更多的在于努力在确定的道路上能够前进多远。而交易者是否能够学有所成，更多的在于努力寻找成功之路的结果。从微观上来讲，有多少交易者，就会有多少交易风格。而这些不同的交易风格，通常并非具有兼容性。也许交易的成功之路就在脚下，但如果你没有经过所有可能的尝试与寻找的话，那么你是不会相信路就在脚下。我曾经就是一位像“贼”一样怀着不劳而获的心态进入股票市场的交易者，而结果真的是血淋淋的。

对于一个商人而言，他需要花费资金来进货，然后再靠以更高的价格销售货物从而获利。也许这个世界上存在一本万利的暴利行业，但绝对不存在无本万利的行业。商人的利润来自于收入减去成本，当商人避免了所有的成本时，也就等于放弃了所有的利润。我的朋友总是玩笑般地问我：“最近又进货了吗？”事实上，我的确是把亏损当做进货而已。我相信，我在整理行情中“垫付”的钞票，一定能够在趋势来临之时连本带利地收回。

有的时候我在回想，小时候炎热的夏天家中并没有空调，全院的居民都会到路边去乘凉。大家轻松愉快地聊着闲天，人与人之间的关系似乎是那么亲密无间。我吃邻居家带来的一块西瓜，他们喝一瓶我家带来的汽水。这种惬意的恬淡生活，似乎永远都与压力无缘。20年后的今天，忙碌的人们在必须努力奋斗才会不被社会抛弃的压力下喘不过气来。同样的夏天，人们各自在家中吹着冷气，似乎都不知道隔壁住着谁。大家在这种令人窒息的压力与忙碌中，变得越来越冷漠……有的时候，我在想，具体哪种生活方式更能够让人的心灵感到幸福与满足呢？经济的发展，追求的到底是什么呢？

请相信我，我们从来都不曾真正地得到什么，我们得到的只不过是靠失去另一些东西而换回来的。交易同样不可能存在不劳而获，更不可能存在无本万利。到底我们用什么换来了真实的利润呢？我告诉你，交易的利润只不过是来自于对心灵折磨的补偿。你可以做一个获利丰厚的交易者，也可以做一个随心所欲的交易者，但你绝对无法做一个既随心所欲又获利丰厚的交易者。

那些为了追求内心满足，而今天“偷一个短”，明天“偷一个短”的交易者，被“剁手”只是早晚的事情。我宁愿将利润隐含起来，也一定把亏损放在看得见的显著位置上。市场有时像一个温顺的小猫，有时又像一只凶猛的老虎。千万别总想着欺负市场，老虎不发威，你还真拿它当病猫呀？

本章的前三节希望告诉大家的是\textbf{：成本观念，以及盈亏同源的观念，这是进入系统交易的门槛，迈不过这个门槛的话，系统交易就是天方夜谭。}对于我来说，一笔交易无论是盈利还是亏损，我做的都是同样一件事情。

\subsection{预测与交易}

{\kaishu
你可以在完全看不懂行情的情况下获取利润，也可能在看对行情的时候亏钱。

\hfill ——金融帝国
}

预测与交易使用的是两种截然不同的思维方式，预测是“点状”的、是相对的，而交易是“环状”的、是绝对的。甚至众多的“交易之环”还要组成更大的“交易之链”。一位交易者是否出色，完全取决于他的“交易之链”是否出色。如果我们不能弄清其中的区别，那么我们就很难完成出色的交易。

对于预测者而言，他们并不需要给出绝对化的二元观点。比如，按照黄金分割，我们可以指出某某价位有阻力，甚至我们还能指出第二、第三阻力位。但阻力位究竟意味着什么呢？预测者并不用回答这样的二元选择，卖出还是持有？而对于交易者而言，必须只能作出“卖出”或者“持有”的二元选择。市场的语言中，并不存在阻力位、超买超卖等词语，而最终必须落实在买卖二字之上。市场不管你怎么想、怎么看，只关心你怎么做！

波浪理论、江恩理论在解释历史行情时，通常都是非常完美的。为了方便探讨，我们只用黄金分割举例。我可以不夸张地说，股指几乎所有的牛熊循环，没有一次是同黄金分割不吻合的。（2245-998）/0.382+998=4262点，与5·30调整前的高点4336点只相差1.7％；（1783-998）/0.236+998=4324，与5·30调整前的4336点只相差0.28％，这不可谓不精确！2245点为上轮长期牛市的顶部，而1783点为大B浪反弹的顶部，998点为大C浪的最终底部，这种分割方式的三个点位选择绝对是极其客观的。黄金分割中的几个比例为0.236、0.382、0.5、0.618、0.764、1这6个数字，而历史行情中符合这几种比例关系的案例比比皆是。这段文字写于5·30调整时，几年后我们得本轮大牛市最终的高点位6124点。有意思的是，我们仍旧能够通过金分割对其进行解释：（2245-998）/0.236+998=6282，与最终的6124点只相差2.6％。

我们假设上证指数最终于1783点见顶，那么按照黄金分割同样可以解释为（2245-998）×0.618+998=1769；（1783-998）×1+998=1783。如果上证指数在2245点见顶，那么按照黄金分割可以解释为（2245-998）×1+998=2245；（1783-998）/0.618+998=2268。如果上证指数在3016点见顶，那么按照黄金分割可以解释为（2245-998）/0.618+998=3016；（1783-998）/0.382+998=3053。

如果是按照切线理论来解释行情的话，那么我们可以画出趋势线、江恩线、X线、平行线等。这么多种切线不可能全部起作用，否则价格将会寸步难行。但通常事后来看，很可能出现一种“切法”能够创造出完美。

事实上，前面章节讲述的众多方法中，更大的比例是预测方法而不是交易方法，毕竟他们还不是“成环”的。交易方法必须是可以绝对化的，而如果一位交易者像分析师那样追求幻觉，那么就有点贻笑大方了。对于交易者而言，观点越绝对越好；而对于预测者而言，观点越模糊越好。

对于预测者而言，他们并不需要“成环”的观点。预测者可以今天看多，明天还看多；而交易者如果今天做多，那么只能考虑什么时候该做空（平多）。也就是说，交易者的观点只能够是多、空、多、空的循环。明白吗？对于交易者而言，看多或者看空是需要资格的！我经常在熊市中听到某股评的观点为看多，下跌半个月后还看多，下跌一个月后仍然看多。我同样经常听到某套牢惨重的股民，对着我说他看多股市。问题是，你还有钱吗？股评任何时候的正确观点都是有意义的，而交易者则不然。很多总是想给我上课的老股民，经常在底部形成后跟我讲：“多么明显的底部特征呀！只不过我没有资金了，而你有这么多钱为什么不去抄底？”虽然我不屑于与他们争论，但我心里很清楚：能够让你看出底部特征的技术，恰巧正是让你现在根本就没有资金的技术。我始终想不明白，总是思考在我的地基上如何盖房子的人是一种什么心态？或许，他们更应该考虑的是，在下次底部来临时，如何才能拥有抄底的资格。

对于预测者而言，他们更在乎正确；对于交易者而言，我们更在乎获利。我经常听到股民说：“因为某某均线金叉，所以我看多。”我觉得这是荒唐的，我作过历史统计，这种单一命题的正确率不足40％。但我确实经常说：“因为某某均线金叉，所以我买入。”这是绝对没有问题的，我虽然没有概率优势，但我有赔率优势。区别在那里？区别在于按照这种方法单次交易是没有任何意义的！

对于预测者而言，他们讲究冲击极限；对于交易者而言，我们更讲究平淡无奇。简单地说，预测者更在乎正确的次数，而交易者更在乎的是赔率的大小。如果利润是一定的话，我希望是由更少的获利次数与更多的亏损次数组合而成。而对于预测者而言，就截然相反。孙子兵法是继承道学文化的产物，其中的一个观点是：他只会去打那些顺理成章就能获胜的战争，而不会去打那种高难度的战争。上涨不停就持多不动，远远没有在横向整理中连续6次指出高低点更吸引观众的眼球。这样一来，预测者更讲究天天有观点（必然是注重短线），而（趋势跟踪）交易者更希望天天没事做。

记得十年前，我在反复阅读江恩理论后，用历史行情进行验证时。我发现历史总是非常完美地符合江恩理论中的观点，如果这种完美能够再现的话，那么非常短的时间内我就能赚取天文数字的财富。有意思的是，其后用江恩理论进行交易时，效果并不理想。更有意思的是，过一段时间后再用江恩理论来解释曾经的行情时，不理想的效果也能变为一种事后的完美。似乎我距离巨大的财富也许只差这么一点点！我拼命地想抓住这近在眼前的巨大财富。

直到发生这样一件事情，才让我深刻地反思其中的奥秘。在上证指数1700点时，波浪理论与黄金分割全部指向完美的顶部。而期货市场又有明显的做空机会，在市场展开短期调整时，我离开股市而进入期货市场。具有讽刺意义的是，期货市场中做空让我卖在地板价上，而反手做多让我两个月内获得了250％的暴利。我开始反思自己是否真的看得懂行情？更具有讽刺意义的是，大盘不但没有见顶反而出现大幅度的飙升。

我突然发现，这种“事后完美、事先不理想”并非一种巧合。按照“树形原理”的观点，我们可以找到众多成功者具备的因素，但你做到这些因素时并不能够保证你能成功。交易的道理也是一样的，你可以找到众多行情形成前具备的规律，但符合这些规律的走势未必能够走出预期的行情。这就解释了，艾略特可以发现著名的波浪理论，但却并不能使用该理论创造巨大财富的原因。历史都是符合逻辑的，但未来是飘忽不定的。《通向金融王国的自由之路》中讲述的事后合理偏向，似乎告诉我们交易者不应该成为历史学家。

不夸张地说，我用了十年的时间，才知晓“预测”与“交易”之间不同的思维方式。十年前我经常在股市的中期底部时满仓套牢，而对明明看出的底部无能为力。进入期货市场以后，我更发现即使看对行情，离市策略、资金管理做不好也照样无法获利。观点的正确与否，只不过是能否获利的一个次要因素。我经历过多次看对行情而无法获利的无奈，后来也经历过看错行情反而获利几十万的辉煌。我经常听见交易者对我惋惜地说：“我曾经非常看好的一只股票变成了黑马，但我却没有拿住。”事实上，对于交易者来说，“做的能力”比“看的能力”更为重要。表面上看，从“看到”到“做到”只不过是一小步，其实这一小步比大多数人想象的要远得多！

大家或许会发现我一个奇怪的观点。市场是值得敬畏的，奢望找到市场的逻辑是愚蠢的。只要你能够形成观点，那么市场就会打败你的观点，除非——你根本就没有观点！“交易之环”是由众多因素组成的：资金管理、开仓、持仓、平仓，并且这些因素是绝对不可分割的。而更多交易者注重的只是对开仓有意义的预测观点。我想说的是，既然交易必须是成环的，那么我们的交易技术也必须是成环的。本章的目的就是希望读者朋友们，能够将此前介绍的交易方法，组合成一种能够成环并且能够反复使用的交易系统。交易的砖块并不能决定交易者的战绩，而由众多技术砖块组成的系统大厦才是真正决定交易者成败的关键。同样的砖块，可以建出风格截然不同的大厦。其实，交易也是一种很个性化的行为，而重要的是你能建出什么样的大厦。

我可以保证，我能用三个月的时间，教会一位不懂交易的新手几乎全部正统的交易技术。但一位新手将这些砖块组建成适合自己居住的系统大厦，就绝对不是短短几年就能实现的。要知道，我只需要三分钟就能教给你一种长期获利的交易系统！同样，让你绝对接受这种“愚蠢”的方法仍旧需要数年的时间。

最后，我还是给大家讲一段我的亲身经历吧。2005年3月我在MACD论坛发表过一篇名为《未来两三年内，股市中可能出现的词汇》的文章。其中指出了未来的市场热点将会围绕着“下游产业蓝筹股”展开；由此产生“深强沪弱”的市场特征；并且深成指能够产生万点行情。别误会，我不是炫耀什么。事实上，我并没有从本轮壮观的大牛市中赚到什么钱，而是在2006年8月就将主战场转移至期货市场。论预测的话，那么我的观点不可谓不经典，要知道2005年3月市场还处在最为恐怖的时期。论交易的话，我根本就没有利用本轮大牛市赚到什么钱。我希望大家能够反思一下，预测与交易之间到底距离有多远？

\subsection{赌博}

{\kaishu
如果你的利益取决于具有不确定性的结果，那么你就是在赌博。

\hfill ——金融帝国
}

我相信，看到我在文章中大谈赌博的读者朋友，一定对我没有什么好的印象。在中国人的传统印象中，“赌博”是一个令人忌讳的词语。也许正是人们对于赌博的深恶痛绝，而使得没有人愿意思考它到底意味着什么。我们还是先来给赌博下一个定义吧：如果你的利益取决于具有不确定性的结果，那么你就是在赌博。所谓不确定性是指：同样的原因可能产生不同的结果。

曾经有人提出中国股市的“赌场论”。事实上这里的“赌”字，只能算是狭义上的赌博。他是指结果完全不具有概率上的偏向性。从广义上来讲，做生意与买彩票都属于赌博。所不同的是，做生意具有正期望收益，而买彩票具有负期望收益。当然有人认为参与具有正期望收益的游戏叫做“博”，而参与具有负期望收益的游戏才叫做“赌”。如果是这样定义的话，那么我是从来不赌的，毕竟我从来没有买过一注彩票。

从广义上看，世界上的所有事物，都无法与“赌博”绝对地脱得了关系。事实难道不是这样吗？否则，我们按照成功者的足迹，自己再走一遍就一定能够成功。这样一来，这个世界上就不会再有穷人了，也不会再有失败者了。相信我，投资收益是不具有再现性的，成功也是不具有再现性的。

肯定有人会反驳我“同样的原因可能产生不同的结果”的说法。古代的哲学家说过：太阳每天都是新的；你无法两次踏进同一条河流；世界上没有两片树叶是相同的……没错，即使是两位先后摇色子的赌徒，他们也不可能是绝对相同。比如，时间是不同的；位置是不同的；摇色子使用的力气是不同的；甚至他们的名字也是不同的。请注意，这些不同之处与赌博结果是没有任何相关性的。这样一来，所谓的“同样的原因”是指任何归纳程度后的“同样的原因”，而不是绝对相同的原因。

有一天上海股市大涨了一百多点，收盘后一位股民说是受股指期货即将推出的传闻影响。我们先不说传闻的真假，如果当日股市暴跌的话，那么可能同样有人说是受这个传闻的影响。通常消息都可以从不同的角度进行解读，而人们认为正确的解读方式又会受到行情结果的影响。在我看来，交易是一门忽略的艺术，而不是一门精确的艺术。我们应该把与结果不具有相关性的原因忽略。否则的话，就会出现：张三赌博盈利的原因在于他叫做“张三”；而李四赌博亏损的原因在于他叫做“李四”。

上面说的这些，是希望让读者朋友们最大限度地理解，为什么我在本书中不止一次谈及赌博这种疯狂行为。下面我们按照职业赌徒的标准，来谈一谈赌博中必须考虑的原则问题。

第一，输不起的不赌。君子算战，先算败。最低限度，我们应该知道可能会出现的最大损失是多少。如果这种可能出现的最大损失是我们无法接受的，那么我建议你不要下注。抱着输不起的心态参加赌博者，通常都会失败。抱着孤注一掷、“不成功便成仁”想法的赌博者，通常也不会有什么好下场。也许一个人的行为，不可能不受到心理状态的影响。如果你真的输不起的话，那么在下注前，你已经输了一大半了。

第二，能够接受任何结果。这要求我们在下注以后，对结果不要有任何期盼。人通常是把应该自己左右的事情听天由命；而又想通过期盼来左右只能由“天”来决定的结果。某心理学家在对拉斯维加斯赌博的赌徒进行研究后发现：通常赌徒都是在下注的时候冲动，而又在轮盘转动后焦虑地期盼。难道我们的交易者不也是这样吗？在开仓的时候冲动，而在持仓的时候焦虑。可能包括我在内的所有交易者，都是将更多的精力放在了没有任何意义的“期盼”上。不是我迷信，你越是在乎的事情，你就越难以做好。

第三，永不赖桌。我们应在在赌博开始以前，就明白自己是在舍什么、是在赌什么。亏损是会永远伴随赌徒一生的现象。亏损并不可耻，可耻的是超过计划以外的亏损。如果亏损已经达到你的计划，这就说明你已经输了！希望捞回一些稻草的滥赌徒，通常会被不可控的风险彻底打败。

第四，赢冲输缩。请相信我，运气是有惯性的！运气好就是运气好，我们没有必要限制自己的好运气。即使好运气结束时会伤害到你的利润，但这时你也是最“输得起”的时刻。这绝对不是迷信，人在不顺的时候喝凉水都能塞牙，而在一帆风顺的时候能够错有错招。甚至做错了事情，还能出现好的结果。不要和命运抗争，如果幸运之神并不眷顾于你，那么你一定要保存实力等到幸运之神发现你的那一天。人应该向乌龟学习，危险的时候就躲在“壳”中忍受一切羞辱，而在捕捉到猎物时，咬住了就不撒嘴！

第五，赌博不应该成为生命的基石。如果人生的基石是建立在不确定性上的话，那么整个人生绝对是一种悲哀。我相信，最能够做得好交易的人，绝对是最能够把交易放下的人。如果一个人把交易本身当做生命的基石，那么很可能你就无法做好交易。人越是在乎、越是依赖的事情，就越难平静地将其做好。能有资格征服、驾驭一种事物的人，绝对是那些能够做到“放下”的人。从人生的角度来讲，人的一生总是把精力浪费在自己根本控制不了的事情上，难道不是一种悲哀吗？

想想看，你是愿意做一个通过概率优势来赌博的交易者呢？还是愿意做一个自以为是在做交易的赌徒呢？

\subsection{再见，确定性}

{\kaishu
如果让我把人生、甚至是整个宇宙浓缩为一个字的话，那么这个字就是“赌”！也许上帝可以不玩骰子，但你我都不能不玩。

\hfill ——金融帝国
}

追求确定性来自于对完美的追求，即使明知这种完美是不现实的，那么也不能让人们放弃追求。更多的人乐意用掩耳盗铃的方式来应对现实世界必然存在的不完美。想想本书开头讲述的那个“手机销售试验”吧，价位、质量、功能、款式是衡量手机的几个标准，他们绝对是互损的。一个方面的优势，必然会出现其他方面的弊端。但消费者宁愿活在店员谎言制造的幻觉中，也不愿意对这种不完美进行取舍。

自古以来，很多人都在争论市场是否能够被预测。一部分人认为市场的波动是杂乱无章，未来绝对是不可知的；另一部分人认为市场的波动是可以预测的，未来是可以计算出来的。但大家是否想过，在可知与不可知之间还有一种状态，那就是，市场的未来虽然不可知，但却具有概率上的规律。我很肯定，市场绝对不是杂乱无章的，同时也并不具有任何绝对规律。

我要说，承认自己看不懂行情并不丢人，甚至在看不懂行情的前提下，实现必然的利润也是很现实的事情。从“赌场原理”可以得知，赌场的经营者根本无法得知与某位特定赌客的赌博结果，但赌场的经营者却知道其能够必然盈利。

很多人奇怪地对我说：“我没有你这么高的交易水平，所以只能做做短线。”简直是胡扯，短线是我都不敢涉及的高难度交易方式。我明白这些人的思维方式，他们误认为我获取的大利润是看得准。他们根本就不相信，我能够在毫不知情的前提下，将获利头寸长期持有。事实上，我的秘密就在于：\textbf{永远用钱去换位置，而不是用位置去换钱}。什么是位置？位置就是在趋势行情中方向上的优势！更多的人，在获取好位置的时候，迅速用珍贵的位置来换取少得可怜、但真实确定的钞票。而我是甘心用一定真实确定的钞票，来换取可能出现的潜在好位置。

我很诚实地告诉很多人，我的大利润不是看出来的，而是赌出来的。有意思的是，几乎所有的人都不约而同地对我所说的这个“赌”字深恶痛绝！我毫不夸张地说，如果让我把人生、甚至是整个宇宙浓缩为一个字的话，那么这个字就是“赌”！也许上帝可以不玩骰子，但你我都不能不玩。当我向一位朋友反复讲述我的理念后，他说出了这么一句话：“你这个不能叫做‘赌’，应该叫做‘博’”。看来，人们对于“赌”这个字的反感程度太惊人了！最终为了避免大家对我的误解，只好将“赌”这个词换一个文雅一些的称呼——“试错”。

曾经的我，为确定性付出了太大的代价。起初我为了每笔交易都必须盈利的确定性，甘心放弃了太多的利润。即使我放弃得再多，确定性也丝毫不会眷顾于我。我用只买入业绩好的股票、只买入大幅下跌后风险释放的股票等等方式，来追求根本就不值得，而最终仍旧追求不到的确定性。后来我使用趋势跟踪策略时，总是等到一条条向好的迹象全部出现后才买入股票。这样做根本无法提高确定性，反而让我的买入价格比系统高了许多。前面提到的“点球原理”告诉我们，寻找确定性最终只能让我们一无所获。在我看来，一个人的交易行为越不需要理由，那么获利就越容易。当我们每多知道一点的时候，就会多付出一些代价。当我们知道全部真相的时候，这种真相已经没有丝毫的意义了。

人们呀，你们追求确定性，追求得真是不计代价、追求得真是毫不理智！追求完美的动机使得你们想绕开所有的缺陷，而不是承担我们必须承担的东西。你可以继续按照“鸵鸟现象”中描述的那样，在弊端中“绕呀绕”。反正我们不乐意相信的东西，就可以当它不存在。说实话，我见过太多的人，用千奇百怪的方法来追求并不存在的确定性。坚强的人们，跟我一起去直面亏损吧！

\subsection{亏损与境界}

{\kaishu
我真的很感激那些连续出现的小亏损。正是这些令人讨厌的小亏损，把那些比我更聪明的交易者挡在了获利的门外。

\hfill ——金融帝国
}

我可以毫不夸张地说：从一位交易者对亏损的认识境界，就完全能够看出他的交易境界。“亏损”是一个那么令交易者讨厌的词汇。但事实上，交易长期稳定获利的秘密就在于对亏损的正确认识。

\textbf{第一层境界：拒绝面对亏损}

生活中有这么一类人，他们永远对自己的错误视而不见，并且总是会在犯错后找出这样、那样的理由来使自己所犯的错误合理化。其实，他们所找出的那些理由是否有道理并不重要，重要的是那些所谓的理由是建立在“自己没有犯错”的前提下。人嘛，就是这样：总是相信自己乐意相信的事情。问题是，这类人反复在这样的心理暗示下，最终会深信自己就是对的。

说到大众对于亏损的态度，让我想起身边的一位朋友。他在一轮大熊市中不断亏损。但有意思的是，与此同时他在事业上却一帆风顺，公司总是时不时地发这样或者那样的奖金。我记得他经常说的一句话就是，“只当”某某奖金公司没有发。言外之意就是，他没有真的亏损，他的损失已经通过另外的渠道找回来了。我相信，这应该不会让读者朋友感到陌生。现实生活中，人总爱在吃亏的时候，为了保护心灵而不断为自己“找心理平衡”。人为了拒绝面对吃亏的现实，总是可以想出一堆千奇百怪的理由：什么破财免灾啦、只当给某某（让他吃亏的）人买药啦……总之，承认亏损是绝对不可能的事情，即使用“免灾”来做平衡也在所不惜。

生活中有什么样的人，市场上就会有什么样的交易者。也许人都是有自尊心的，似乎每位交易者的成长道路中都会经历这样的阶段。同样，他们对待市场上的原因根本无法客观地分析。简单地说，他们不是因为看涨而做多，而是因为做多而寻找看涨理由。这样的交易者在不断对自己心理暗示的催眠下，慢慢地使自己的信念同市场的客观现实越走越远。如果市场最终的走势真的和他们所希望的结果一致，那么对于他们也不会有任何好处。这样会使他们的自尊心膨胀，下次就更难面对自己的亏损。想和市场一争高下的动机，早晚会被值得敬畏的市场所扫地出门。没错，这只是时间的问题。

\textbf{第二层境界：努力规避亏损}

这类人在生活中应该算得上“完美主义者”，他们根本不允许自己犯任何错误。如果真的犯了什么错误，他们就会陷入反复的自责中无法自拔。也许他们的目标就是努力使自己在未来能够不再犯错。

市场上的这类人已经比第一类人进步很多了。他们已经明白在市场上的生存法则是通过敬畏市场使自己的交易能够保持弹性。他们明白市场永远是对的，和市场“理论”是一件愚蠢的事情。但他们还是没有跳出“对错与盈亏具有完全相关性”的怪圈。他们对自己过于苛刻，而将自己根本无法左右的事情，让自己来负责。如果说第一类人是在和市场较劲，那么第二类人就是在和自己较劲。他们认为亏损是可耻的、是应该避免的，或者说自己应该为亏损负全部责任。这类交易者可能会陷入自责的痛苦中而无法自拔，并且最终成为“学习效应”的受害者。甚至，他们会为了规避1元的损失而心甘情愿地放弃3元的利润。很显然，接受亏损的同时并且享受利润，应该是更为划算的选择。

\textbf{第三层境界：承认亏损的合理性}

如果说生活中的完美主义者是那些把避免犯错看得比获得辉煌更重要的人，那么做起事来就会畏首畏尾。什么都不做，就会什么都不会做错，但同样什么也都不会做对。或者说，即使做事情，也绝无胆量去做任何远离羊群的事。事实上，这类人更热衷于去做大众都在做的事情，并且力争能够做的更好一些罢了。

我在中专时有位好朋友，他的能力非常强。现在他有自己的公司、房子、汽车，要知道他也才25岁。当年朋友们对他的评价是：拿碰壁当“找乐”。也许正是因为他不忌讳碰壁、不忌讳犯错，所以乐观的“冲劲”才使他能够打拼出属于自己的成就。我和这位朋友有很多的相似之处，比如：我们的家境都不优越，我们的学习成绩都不好……有意思的是，当年上学时家境优越、学习成绩优良、甚至充满优越感的同学们，就基本上找不出什么有所作为的人。我想这其中的奥秘或许就是，这类人从小到大就从来没有输过，所以自尊心强得要命。或许他们有更高的起点，甚至学会了所有的技能，但惟独最重要的一件事被他们所遗漏：那就是学会如何去输！与此对应的是，我们的成长轨迹似乎让我们充满了忧患意识，“要想不淹死、就必须往火里跳”则成为了我们的人生格言。事实上，我们更加自信，但我们的自信从来就不是依赖于总是获胜。我们早就习惯了超脱成败的结果，来保留我们的自尊与自信。

在市场上，如果你能承认亏损的合理性，并且能够接受亏损是无法避免地会发生。那么你的交易行为才不会畏首畏尾，从而能够做到得心应手。否则前怕狼、后怕虎的拖拉作风，只会让我们最终等到真正的狼、真正的虎。但无论如何，这时的交易者还是希望最大程度地规避亏损。只不过他们把获得利润看得比避免亏损更重要。

\textbf{第四层境界：认识到盈亏同源}

这时的交易者已经意识到盈亏是我们左右不了的事情，我们只能从宏观上来寻找所谓的概率优势，而在微观上的得失只能交给“天意”来决定。无论是盈利还是亏损，我们都是在不断地做同一件事情。当我们规避了亏损，同时也就规避了利润。如果说我们的交易方法是具有正期望收益的话，那么避免某次亏损就是一件并不划算的事情。也许这时的交易者，已经开始思考交易的本质到底是什么？人类的智慧在市场面前到底意味着什么？

\textbf{第五层境界：得失随缘、心无增减}

当交易者已经明白，除了寻找宏观的概率优势以外，我们对市场根本无能为力。如果说这时的交易者已经能够“爱上你的亏损”，可能显得有些夸张。但最低限度，应该可以做到“得失随缘、心无增减”。这已经到了我所知道的交易最高境界：放下交易，而跳出市场之外。这是一个心灵成熟的标志，需要很长时间来慢慢磨炼。对于交易的结果而言，交易者已经从“依赖”转变为“顺其自然”。对于交易信仰而言，交易者已经从与市场的“对抗与争斗”转变为“包容与融合”。这就是道学文化中所谓天人合一的最高境界！

\subsection{道与术}

{\kaishu
一流的交易策略搭配上二流的交易技术，要比一流的交易技术搭配上二流的交易策略好得多。

\hfill ——大意来源于《克罗谈投资策略》
}

自从看完《通向金融王国的自由之路》后，我开始大量阅读国外书籍。我发现国外的书籍与国内的书籍有很多本质上的不同。如果说“道”与“术”最早出现在中国古典哲学中的话，那么真正从中受益的却是外国人。我追根溯源后才恍然大悟，原来这些宝贵的思想最早是出自老子的《道德经》。

“所谓天才，就是99％的汗水加上1％的灵感。但往往起决定性因素的不是99％的汗水，而是那1％的灵感。”我敢肯定，这句爱迪生的名言，至今仍有超过一半的人不知道后半句。或许对于交易而言就是，99％的技术加上1％的理念，而决定你最终成败的并非是99％的技术，而是那1％的理念。对于国内交易者而言，往往并不缺乏刻苦的努力，真正缺乏的是对努力方向的思考。

往往越接近本质的东西，就应该是越简单的。无论是交易者还是科学家，比拼的都应该是化繁为简、寻找本质的能力。宇宙的确是纷繁复杂的，但支配这种复杂的却是屈指可数的四种作用力。爱因斯坦还是觉得用四种作用力来解释宇宙现象仍旧不是本质，因此后半生的精力都放在研究“大统一理论”上。如果我们回顾物理学的发展史，的确可以发现这种化繁为简、寻找本质的轨迹。早在1687年，牛顿就将苹果落地与天体运行的作用力归纳为一种叫做万有引力的东西；在19世纪中叶，电和磁还被看成是两种独立的事物，而麦克斯韦将他们统一为电磁力；上个世纪60年代，三位科学家提出了弱电统一理论，从而将电磁力与弱相互作用力统一起来。科学界有一个奥姆剃刀原理，它告诉我们简单的才是最好的。

与此相对应的是，交易者大多都有将简单的事情搞复杂的倾向。如果不够复杂的东西，给人的第一感觉就是不深奥。我相信能够看到这本不算热门的交易类书籍的朋友，应该都是对交易研究得非常刻苦的人。那么我来问一个问题：你认为自己能够在市场上获利的基石是什么？我相信，很多交易者都没有认真地思考过这个问题。大家研究的只是一种种相互独立的交易技术，而从来没有关于这个问题的核心逻辑。如果你为了回答这个问题也要写一本书的话，那么应该可以说明你并没有找到交易的真谛。我问过身边朋友类似的问题，但我从来没有听到过相对满意的答案。有的答案是关于基本分析的，挑选有潜力的上市公司……但对于如何挑选有潜力的上市公司，他恐怕无法给出一个一致性的标准。就算我要求听听一致性的标准，那么这种标准也不是一两句话能够说得清楚的，甚至复杂到我都无法验证其有效性的地步。

如果让我回答这个问题的话，我的答案将会是非常简单的：亏损兑现、盈利挂起！我没有孔老夫子的水平，能够将复杂的道理归纳为一个词、甚至一个字。但愚钝的我也可以用八个字来归纳本人交易理念的灵魂！而更多交易者的理念可以用庞大而没有灵魂来形容。再复杂的科学理论，往往都只不过需要几块屈指可数的基石（假设）。如果作为基石的假设是成立的，那么整个理论大厦就是坚不可摧的。当然，撼动整座理论大厦也只不过是需要证伪其中一个假设罢了。我很难相信，大部分交易者的理念大厦往往是在没有基石的情况下建造起来了。这些交易方法的有效性从来无法被证伪，因为从一开始同样是无法被证实的。

当然，如果让我说明白我的基石为什么能够获利，恐怕就需要多说几句话了，这个问题我会留到后面作出系统的回答。其实，我使用的交易技术仍旧是非常简单的，只不过是如何量化这八个字罢了。我不想在交易上追求极致，往往极致的东西都不会是稳定的，甚至隐含了巨大的隐患。所以即使是对于交易技术而言，我也坚决地执行“奥姆剃刀原理”。能够用简单的方法来解决，就绝不使用复杂的方法来解决。我曾经看到《克罗谈投资策略》中的一种观点：一流的交易策略搭配上二流的交易技术，要比一流的交易技术搭配上二流的交易策略好得多。而我的观点更加极端：技术与策略同样是互损的！如果你具备了一流的交易技术，那么你就注定与一流的交易策略无缘。而如果当你真的掌握了一流的交易策略，那么你根本就不可能容得下一流的交易技术。为了形象地说明这其中的奥秘，我还是给大家讲一件我曾经工作中的事情吧。

前面已经说过，我的第一份工作是做手机行业的销售代表。其实我的销售理念很简单，就是把每一款新上市的机型拿给区域内的经销商看一看。客户进货就进货，不进货也没有关系，总之我不会死皮赖脸地软磨硬泡。即使客户不进货，我仍旧按时地解决好售后问题；对于进过货的客户呢，公司有专柜（柜台）或者其他政策的时候，我都是无条件地让客户享受。当然，如果是其他的销售代表，一定会利用公司的政策作为筹码，诱使经销商进更多的货。很显然，一开始我完成的销量一定会比同事小得多。但渐渐地，由于我“实在”的风格，几次合作后，经销商对我越来越认可，觉得同我合作非常放心，而我做的品牌也渐渐在市场上卖火了。其实，我能够取得经销商的信任，只是因为一件很巧合的小事。记得那时有一款性价比不错的机型，但这款机型存在一个高返修率的隐患。在客户进货时，我将这个问题说得很清楚，并且建议不要压货太多。虽然只是少量进货的试卖，但返修率已经将经销商弄得焦头烂额了。或许，正是我有悖于其他业务员的做法，让经销商非常感动。而正是这件小事，让我明白了“做渠道”与“搞推销”之间本质的不同，而更多的业务员则是用搞推销的模式来做渠道。虽然我那个“少进货”的建议让我少赚了一些工资，但在以后几年内我所得到的回报简直大到了不成比例。

随着销量的增长，这时出现了一个新问题，经销商之间的竞争往往会引发“砸价”的问题。他们通常会将自己不压货的机型报出低价。即使没有利润也没关系，反正他们也不是真的想卖，而会说这款手机质量不好什么的，然后再主推他们库存量大的主卖机型。我核心的经销商也和我探讨过这个问题：如果我的品牌被砸得利润很小，那么店员就没有推销的热情。

我认真地思考了一下这个问题：诺基亚可以用全线开花的销售策略，毕竟很多顾客进店后就点名要这个品牌，至少没有经销商敢说诺基亚质量不好。但国产品牌不行，必须靠人来推销。事实上，很多经销商都成为国产品牌的“地包商”，其他商家都是从“地包商”那里分货。很显然，经销商一定会不遗余力地推销自己包销的品牌，而其他品牌就会沦为“砸价”的工具。如果手机没有利润，那么店员就拿不到多少利润分成。或许，我可以每天跑到卖场给店员买些冷饮讨好他们；然后在各个经销商之间周旋，尽量避免“砸价”。其实我很清楚，任何市场的参与者都是求利的。就算我做再多的事情，也只不过是表面上有所改变。而只有更高的利润，才能够转换成为店员推销的热情。最终我作了一决定：背着公司将进货门槛提到一个很高的水平，这样就能够把不是真心想卖货的经销商挡在门外。

这样一来，我负责的市场渐渐变为一家经销商“独卖”，而利润对于各个环节来说都算是比较丰厚了。店员由于利润的驱动，从“要我卖”变成了“我要卖”。现在回想起来，那段时光仍旧让我津津乐道。我每天几乎什么事情都不用做，甚至经销商自己去公司提货、自己派人去跑售后。那么我在干什么呢？哈哈，我在家里炒股票呀。当然，还有就是和经销商吃吃饭、蹦蹦迪、打打牌什么的，而且我们在娱乐的时候从来不聊工作上的事情。我创造了一个同事们不敢想象的奇迹：那就是别人白天“下市场”，而我是晚上“下市场”。道理很简单，白天经销商有很多事情要忙，哪有时间陪我玩呀。大家或许会觉得不可思议，在绝对的买方市场中，怎么会出现买家自己提货、跑售后的事情呢？道理其实很简单，经销商既然开店就一定是求利的。而我如果继续求利的话，那么一定会存在冲突。既然经销商求利，我就必须“错开”需求，转而追求少干活。其实，经销商也有专门的员工每天跑市区，顺便帮我跑一下并不增加成本。而对于我来说呢，把除了工资之外的所有潜在收入都贴补给经销商。这样一来，我自然可以挤出时间在家炒股票，并且最终实现与客户的双赢。这里解释一下，所谓的潜在收入大多都是摆不到台面上的，但几乎所有的业务员都会钻这些空子赚点小钱。事实上，倒未必是经销商在乎我这点小钱，而是我的风格让经销商感到非常舒服、甚至是感动。毕竟，经销商请我吃吃喝喝花掉的钱，要比这些所谓的潜在收入多得多。有意思的是，在我极其懒惰的工作表现下，我的销量却反而越来越大。甚至公司领导都非常认可我的销售策略，不止一次跟我的客户承诺不会向其他经销商供货。

事实上，对于渠道类的销售工作来说，的确存在二八法则的效应。你付出20％的努力，就能够产生80％的销量。而如果你想要得到另外20％的销量，那么你就要继续付出80％的努力。客观地说，我的绝对销量在公司中只能算得上中等偏上的水平。如果是考虑当地市场占有率的话，那么我的排名就在公司名列前茅了。这是因为我负责的是一块很小的市场。当然，如果是大一些的市场，我的“独卖”策略也不可能增加太多的销量。至少我的销售策略非常适合我所在的市场特征，毕竟那是一块“手机品牌数量远远多于经销商数量”的小市场。如果我只是想教大家如何在工作上偷懒，那么我就真是够“没流儿”的啦。其实，后面发生的事情更加精彩。

或许是因为分公司总经理的私心，那时公司几乎不计代价地追求销量，而从来不考虑销量到底是良性的还是恶性的。为了增加销量，公司简直无所不用其极。公司派人装成顾客，到没有进货的商家处询问我们品牌的手机，从而使得经销商错估市场潜在的销量。这种策略只能够解决短期问题，而副作用就是会出现货物的积压。作为老业务的我很清楚这一点，曾经发生过很多次只因为“历史遗留问题”无法解决，而无法继续合作的案例。哪怕只是多年前积压下来的一两台手机没有销售出去，只要公司不给退货，经销商是不可能再继续进货的（我们公司又没有退货的机制）。退一万步说，就算手机被顺利销售出去。但只要花费的时间比经销商预计的更长，那么也会影响第二次进货的热情。

我的销售理念是，当“渠道”和“销量”发生冲突的时候，永远是舍销量、保渠道。如果我明知道某个型号的手机不好卖，那么即使经销商主动要货，我也会劝他再考虑考虑。因为我很清楚，渠道内的货如果卖不出去，将来还是给我自己“找病”。这与其他销售代表不计代价地拼命压货形成了鲜明的对比。这就是我追求自然销量的原因，而我所能做的就是不会因为渠道的问题而影响到自然销售而已。如果我非要违背自然销量，而刻意努力提高销量的话，恐怕即使在终端销售环节下手，也得不到预期的效果。最多也只不过是我在场的时候，店员“局着”面子帮我推销一下罢了。何况我又没有分身之术，不可能同时出现在另外几家分店的卖场。我从不奢望把梳子卖给和尚，这种努力是没有任何意义的。在我无法改变根本结构的前提下，硬要改变结果同样是没有任何意义的。

有意思的是，那时几乎是我们这个品牌，在手机市场中最辉煌的一段时间。其他区域的经销商看到潜在的利润，也希望成为我们品牌的“地包商”。公司的做法更是离谱，居然做出了“一个闺女，许给两婆家”的事来。你想当“地包商”没有问题，但需要进多少多少台货。还有人想当“地包商”，只要从公司大量进货也没有问题。这时我已经很清楚了，这块市场已经彻底完蛋了！虽然囤货的商家对公司非常不满，但那时我们品牌的确挺好卖的。看在钱的面子上，虽然总是向公司发泄不满，但谁也没有和公司撕破脸。

我的印象非常深，开大会的时候，总经理拼命地表扬那个区域的销售代表。与此同时，公司很看不惯我的销售思路，甚至没有点名地在会上说：业务员要以公司的利益为重，不能胳膊肘往外拐。我能怎么办呢？明知道听公司的就是死路一条，而总经理应该不会愚蠢到不知道后果，所以只能假装听不懂，并且顶着压力继续按照我的思路来运作市场。最后公司实在受不了我了，准备把我炒掉。在我不知情的情况下，公司派业务经理走访我的市场。与我的客户讲明来意以后，我的客户直接说：“你就把他炒掉吧，反正我以后也不会再做你们的品牌了。”这位经销商早已成为我的好朋友，他的做法并不让我出乎意料。让我出乎意料的是，经理到另一家早已不与我们合作的经销商处，那家经销商仍旧对我评价非常高！要知道，我曾经跟这家经销商因不给他放货而发生过争执。很显然，公司立刻打消了炒掉我的念头。

事后我问我的经销商朋友：“我们经理找你谈的政策肯定是我所给不了你的。即使没有政策，现在卖我们的手机都很赚钱，你为什么如此生硬地拒绝了呢？”我的经销商朋友是这么回答我的：咱们先抛开感情不谈。我相信再来一位业务员会对我很好，甚至比你对我还要好。但问题是，他在这块市场立住了以后，一定会为了追求销量而不规矩地把货放到其他经销商那里。咱们费了这么多年的时间才把这个品牌炒起来，“钱道”都没了，我还要钱干什么呢？其实另外一家经销商对我评价很高的原因也很简单。我在解决他的售后问题时，如果能够变通一下不收费的，就免费帮他修一下。而如果是其他业务员，一定会找公司免费修完以后再找他要钱。何况，他有自己包销的品牌，不可能大量从我们公司进货。当时的这种模式应该是对他最有利的。毕竟他不会傻到花大笔资金进货，而只是为了和我的经销商斗气。事后我还在思考：如果业务经理小家子气一点，只要求他进三五台货，那么事情的发展会不会出现另一个版本呢？毕竟，他早就对我的经销商独卖赚取高额利润，而我又不给他放货不满。但当进货的权利与销量的义务挂钩以后，他的确作出了明智的选择。我真的感叹道，这就是与强者合作的乐趣。你只要站在他们的角度考虑问题，他们绝对能够算得清楚自己的利益。

后面的事情，几乎和我预料的丝毫不差。花无百日红，公司产品线的竞争力开始下降。我的同事开始变得焦头烂额，那些曾经为了利润而忍气吞声的经销商集体爆发了。别说继续进货了，就是消化掉以前的库存都成为很大的难题。那位“以公司利益为重”的销售代表，并没有得到公司的谅解，不得不离开公司，而我们的总经理却在销量最辉煌之时调到了其他城市的分公司。由于我的渠道是完全良性的，从而继续着以前的模式。几年以后，随着手机行业的没落，分公司最终被撤消了，我也就失业了。由于我与经销商朋友多年的交情，他又帮我在其他品牌的代理商处找了份工作。总之，我的第一份工作在干了五年以后还是得到善终了。毫无疑问，我当仁不让地成为了公司最老的一位销售代表。这五年来，作为N朝元老的我，看着一波又一波销售代表的来来走走，而我却成为了公司的一个传说。

听了我这段故事以后，大家是否能够理解道与术的区别呢？或许，你再多付出80％的努力，的确可以得到那剩下的20％销量。但这最后20％销量的获得，真的没有任何隐患吗？作为销售代表的我，本质工作只是维护渠道，剩下的事情我从来都是听天由命。或许，正是我无为的风格，才让我最终得到了善终。如果说因为我没有维护好渠道，从而使得可以销售良好的手机无法销售的话，那就是我工作上的失职。但在手机没有竞争力的时候，我“有为”地与市场对抗也没有任何的意义。不如在保证渠道良性的情况下，等待下一次品牌竞争力的春天。

客观地说，当产品线竞争力下降的时候，虽然我的绝对销量也会下降，但我的相对（其他同事的）销量却会增加。道理很简单，在产品好买的时候，更低的利润在容易推销成功的衬托下并不是一件非常难以接受的事情。而当产品不好卖的时候，就只能靠高利润来实现销售了。这很类似于开放式基金的销售模式，我们似乎找不到任何开放式基金相对于封闭式基金的优势。封基具有更低的交易成本、更快捷的交易模式，甚至还有不小的价格贴水空间。我们很容易想象，如果开放式基金的交易成本降低到封闭式基金的水平，那么销售开放式基金几乎就成为不可能的事情。而现实是，开放式基金的推销者可以拿到交易成本80％到100％的佣金收入。正是这块丰厚的收入，才使得开放式基金的销售成为可能。发现没有？靠人来推销的产品，一定是最贵的！另一个方面，经销商为了捍卫自己在当地市场“独卖”的地位，也必须在产品线没有竞争力的时候维持一个说得过去的基本销量。当然，这必须建立在绝对信任的基础之上，毕竟很少有人能够在销售火爆时兑现曾经的承诺。甚至就连分公司老总的承诺，在短期利益面前也会变得一钱不值！事实上，除了利润以外，推销我的品牌还有一个更加明显的优势。只要客户真是有心想买，那么就算再去其他商家处询价，最终仍旧只能回到曾经给他耐心讲解的店员处购买。与此相反的是，店员耐心讲解其他品牌的手机，最终很可能是在为别人做嫁衣。这或许说明，当你为了实现某种目的的时候，一定要舍弃一些现有的东西。而问题的关键只在于，舍弃的与得到的哪个更有价值？

凭我同经销商多年的交情，即使很难卖的手机，只要我开口让他进个几十台货肯定没有问题。但我在生意上从来不“卖脸儿”，因为这只能解决一两次问题，而不能够改变经销商需要追求利润、店员需要赚工资养家的本质。明明人家卖别的手机可以赚更多的钱，那么我与市场法则对抗真的有意义吗？实现良好销售的本质只在于销售者得到的高利润，而通过其他方法追求销量只会让你离本质越来越远。

大家想想看，我在工作上的成功是因为我多做了什么吗？不，完全不是！而是因为我懂得应该少做些什么！事实上，我只不过做到了诚信，以及不是什么钱都赚。大家反过来想想：如果你什么都不干，是不是一定就是诚信、以及不是什么钱都赚呢？如果非要说我多做了什么，那恐怕就是不给其他经销商放货。甚至就连唯一的经销商进货时，我有时都会建议他“少进货”。问题是，我的所作所为，似乎没有一条不是有损微观销量的！很多同事的努力表面上看是在追求更多的销量，而实际上却是在不知不觉中把自己的路全部堵死。当他们走向辉煌的同时，也就注定了他们会很快离开。

我通过经销商朋友的介绍，找到第二份工作时，同事们似乎没有人相信，我是一位化石级的老业务员，因为我似乎对于销售技巧一概不知！或许几年前我刚工作的时候，也学习过如何讨好店员、忽悠老板、在谎言中周旋的本领。但真正“得道”的我，早就把这些小儿科的雕虫小技忘得一干二净了。与此类似的是，在火爆的大牛市中，朋友聚会的饭局总会谈到如何炒股赚钱的问题。总之，他们说的头头是道，我却根本插不上嘴。我经常说，咱们聊点别的吧，讨论这个没有意义。为了消除无聊的尴尬，我还和他们开过这样的玩笑呢。在他们讨论到创投概念股时，我半开玩笑地问道：“什么是创投呀？是磕脑袋吗？”事实上，前半句话还真不是玩笑，那时我还真的不知道创投是什么意思。其实，我并非是不想帮助身边的朋友靠炒股赚些钱，而是我对于改变他们的核心思想根本就无能为力。他们所追求的东西与我有本质的不同，既然问题都是不同的，那么探讨答案还有意义吗？

在我看来，道与术的区别在于：一个是求（钱）道而舍钱；另一个是求钱而舍（钱）道。在工作上，我求渠道而舍销量；在市场上，我求位置而舍钞票。因为所追求的东西不同，所以自然一个是简单的，另一个是复杂的。难道不是吗？在工作上我没有办法将我的风格与他人的风格兼容起来从而如虎添翼。同样的道理，在交易上也不行！当你用对抗之心寻找技术的时候，你找到的技术一定是复杂的；当你用顺应之心寻找技术的时候，你找到的技术一定是简单的。如果你为了达到某种目的而不得不持续付出巨大的努力时，那么你的模式一定是没有顺应事物的本质。

在我看来，交易者的问题不是他们在成功之路上跑得不够快，而是从根本的方向上就已经错了。如果说一个人的交易理念是正确的话，那么即使采用最粗糙的武器也能获得成功；而如果一个人的交易理念是错误的话，那么他使用再精良的武器也摆脱不了最终失败的命运。那么交易之道是什么呢？如果你还不明白，那只能说明你的悟性太差了。我已经说得够通俗了，是吧？我真诚地希望更多的交易者，不要在那些细枝末节上继续浪费时间，不要再去做那些我们根本就不可能做到的事情。

\subsection{风险防范}

{\kaishu
我的风险防范策略只有两点，一个是不丢掉自己的位置，另一个是绝对的止损点。

\hfill ——金融帝国
}

风险防范对于任何交易者而言，都是一个最为重要的问题。对于大多数交易者而言，他们会用“传统”的方式来进行风险防范。那么传统的风险防范方式到底能够起到一个什么样的效果呢？

\textbf{分散投资}

分散投资可以作为资金波动幅度的“减震器”，从而规避大面积的资金回撤，并使投资收益实现较好的稳定性。如果你的操作方法期望收益是一定的话，那么分散投资并不能起到规避风险（提高期望收益）的作用。道理很简单，如果你（的方法）是错的，那么十个你也是错的。只不过分散投资可以使你更稳定地亏损，而不会使得最终结果出现与理论值太大的偏离程度。也就是说，分散投资是一种资金管理工具，而不是一种风险防范工具。

我见过很多股民都是“盈利兑现，亏损挂起”，但他们又通过分散投资来规避市场可能会给他们带来的致命打击。这样一来，这类交易者就买入十只以上的股票，哪个获利了就兑现，而亏损的就不去管它。在我看来，这种可怜的分散投资只能是让交易者更加平稳地亏损。更加可怕的是，这种策略在不断过滤下，注定会使你的持股变成弱势股的集中营。

\textbf{买入优质股票}

无论你相信与否，现在的每股收益与未来的涨跌并不具有任何的相关性。如果非要把股票的涨跌与基本面扯上关系的话，那么未来的涨跌也只和业绩（预期）的变化呈现正相关。如果市场上真的有某只业绩优良的股票以不可思议的低价出售，那么我通常认为这是一场业绩的骗局。事实上，很多新股民在对每股收益的崇拜下，都会陷入这种业绩骗局。

\textbf{买入便宜的股票}

所谓便宜有三种含义：买入价格低的股票；买入下跌后的股票；买入疲弱的股票。我很不理解股民为什么会有这样的误区，很多人告诉我高价股亏损十几元很轻松，而低价就算“跌没了”才几块钱。难道这些股民炒股票不看幅度而看绝对的价格吗？总体上来看，低价股的想象空间同风险都要高于高价股。在《短线交易大师：工具与策略》中有这样一个故事。一位交易者给作者打电话时指出：“我并不否认你们推荐股票的正确性。问题是这对我没有任何的好处，因为他们通常都是正在上涨的股票。”事实上，买入正在下跌股票的这种癖好绝对是一个误区。从概率上来看，其效果绝对不如买入正在上涨的股票。这是一个经典的案例：交易者相对于获利而言，更在乎自己信念体系下的“正确”。买入弱势股而占市场的便宜更是荒唐，要知道市场上的强弱互换可能需要五年的时间。

上面的内容前面都有所提及，这里就让大家回顾一下吧。既然这些不能实现风险防范，反而成为交易者加速亏损的思维误区。那么什么样的方法能够防范风险呢？在我看来，所谓风险是来源于两个方面：一个是错误的买入，一个是错误的持有。规避错误的买入又有两个办法，一个是提高买入的正确率，另一个是减少买入次数。如果说赔率与正确率本身就是互损的，那么就只剩下减少操作次数了。避免错误的持有，就是永远不让自己被动，那么唯一可行的方案就是绝对的止损点。

答案很简单，我的风险防范策略只有两点，一个是不丢掉自己的位置，另一个是绝对的止损点。不丢掉自己的位置从某种意义上来讲，就是最好的风险防范策略，而绝对不只是获利的工具。对于2007年强劲的股市而言，很多股民通过“截断利润”来规避风险。当然如果你卖出后就与股市说再见，那么你的风险防范策略还是可行的。问题是如果卖出后还要买入的话，那么你仍旧会面临同样的风险。而你买入的次数越多，那么买入在最高点的概率就越大。

面对4400点的股市，对于丢掉位置的我，会感到非常迷茫。如果现在买入，那么有可能会买在最高点。如果不买入的话，那么大盘仍旧可能不断创出新高。最好的风险防范就是1800点买入股票而一直持有到现在。这绝对不是事后诸葛亮，如果我们绝对不丢掉自己位置的话，那么这并不难做到。长赌必输，你交易的次数越多，就越有可能做出糟糕的交易。而如果你对糟糕的交易又没有任何“免疫力”的话，那么你一定会输得很惨。保住位置的优势在于：什么都不做就什么都不会做错，并且什么都不做利润也会不断地增长。有趣的是，犯错误也同样是需要资格的。当你保证不丢掉自己在趋势中的位置时，你也就失去了犯错误的资格。

\subsection{套利原理}

{\kaishu
套利思维与交易系统的不同之处在于，一个是同时点上的盈亏对冲；一个是不同时点上的盈亏对冲。

\hfill ——金融帝国
}

在学习期货从业资格考试的过程中，我必须学习很多关于套利的知识。套利的方法很多，包括跨期套利、蝶式套利、跨市套利、跨商品套利、期权套利等。套利可以设计出很多的模型，似乎对于交易而言只有套利完全具有科学的色彩。在学习的过程中，我发现套利思维有一个明显的可取之处。

能够影响市场未来走势的因素无穷无尽，而我们又不可能完全了解和运用这些无穷的因素。套利思维最大的优势在于：它能够用对冲的方式来化解大量的不确定因素，从而利用某一特定具有优势的因素在市场中取利。比如，跨期套利就是在回避能够左右未来涨跌的大量因素后，利用不同交割月间差价不合理的因素来获取利润。这充分体现了忽略的艺术，从而使交易者能够在完全看不懂未来走势的情况下，利用某种优势来从市场中取利。与此有异曲同工之妙的就是咬死策略，它同样能够回避市场上的大量不确定性，而只从有概率优势的角度入手。比如，指数化基金就是根据经济长期必然增长，从而使得股票指数长期而言具有涨跌不对称性的概率优势入手。至于选股、中短期走势等不确定性因素完全采取“无为”的忽略艺术。从国外的基金业来看，能够跑赢指数的投资基金少于50％。并且随着运作时间的增长，能够跑赢指数的比例会呈现越来越少的趋势。我突然想起了郭德纲的一段相声：“有什么不会的问题尽管来问我，我会拣着我会的问题来回答。”其实，交易也是别无二致的，将我们回答不了的问题对冲掉，同时对我们能够回答的问题下注。

但传统的套利思维又有一个明显的缺陷，那就是必须利用市场的不合理因素来实现获利。这意味着套利者是在对市场指手画脚，告诉市场应该如何去做。这种自作聪明而没有弹性的交易思维必然会引向两种结果。第一种结果就是无风险的利润少得可以忽略不计，并且交易行为繁琐。第二种结果就是利润有限而风险无限，也就是说市场不合理的程度是有限的，而市场继续扩大不合理的程度是无限的。也许完美的套利模型真的只能存在于书本上。自作聪明地找出市场的不合理，与我敬畏市场、随市而动的交易风格格格不入。所以我个人是从来不会采用传统意义上的套利交易。套利是不可能具有弹性的，你因为市场的不合理而开始套利，那么你应该不会因为市场的更不合理而结束套利行为。也许套利的结果就是能够产生很大的正确率，但无法回避被市场扫地出门的风险。要知道当猪吃的玉米比人吃的小麦还贵的时候，那些利用玉麦比价不合理进行套利的交易者必然死无葬身之地。

其实套利行为还有一个特点就是盈亏对冲是同时点的，你在一只脚出现亏损的同时，你的另一只脚能够立刻对冲掉你的风险。如果将套利思维稍微加以变形，就能够成为我们通常所说的交易系统。也就是说，当你把利用市场微观上的不合理性转变为市场宏观上的概率优势；同时你再把同时点的盈亏对冲转变为不同时点的盈亏对冲。这时的套利思维和交易系统就别无二致了。

在市场上确定性是相对的，而不确定性是绝对的。套利原理的精髓就是在于不是与不确定性对抗，而是在最大程度上将不确定性包容。其实一位交易者的成熟过程基本上能够体现出这样的轨迹。从对不确定性的逃避到与不确性的对抗；从与不确性的对抗到对不确定性的接受；从对不确定性的接受到对不确定性的包容。最终交易者能够超越不确定性的问题，而利用哪怕是一丝的概率优势来实现长期稳定获利。这时的交易者已经不是非要利用必须看懂市场不可的预测来实现获利。

\subsection{一致性}

{\kaishu
即使你强硬地约束自己的行为，你那丑陋的本能仍然会不时地出来找麻烦。要知道越复杂、越相对的约束规则，就越难以真正地控制自己的行为。

\hfill ——金融帝国
}

人的心灵是一部世界上最为复杂的机器。对于一个人而言，对外我们无法全部地认识宇宙；对内我们无法全部地认清自己。也许人的一生就是在不断地做着两件事情：一件事情是对外界的探索，另一件事情就是对真我的寻找。我们应该永远是一只眼睛看着世界，而另一只眼就看着自己的内心。对于我们的心灵而言，唯一可以肯定的是：我们不知道的远比我们知道的多得多。

人性是市场的一面镜子，同时市场又是捕捉人性的天敌。如果说人性是无限复杂的话，那么市场也应该无限复杂。当然我们能够找到人性的规律，自然也就可以找到市场的规律。问题是，我们无法挖掘出市场全部的规律。在值得敬畏的市场面前，我反复地发现“墨菲法则”的身影。对于我现在的境界而言，我不得不把人类思考的价值定义为负值。也许我说不清楚原理，但我肯定市场总是在捕捉我尚未意识到的本能偏向。

前面我已经说过，市场上20％的股票要大大强于大盘，而80％的股票要弱于大盘。让我感到惊讶的是：身边的股民通常没有人持有大大强于大盘的股票。我在反复探索其中的奥秘，并且也发现了其中大部分的原理。我可以轻松地设计出选股程序，并且建立自己的股票池。股票池中股票的平均表现，通常都会比全部股票的平均表现要强劲很多。但这种选股程序只能选择出股票池，而绝对不可能指向最后实施买入的股票。这时我发现了一个奇怪的现象：当我从股票池中挑选买入的股票时，挑选的结果仍旧会落后于整个股票池的平均水平，甚至还可能出现落后于全部股票平均水平的情况。

这是一种惊人的现象！实际操作的结果，通常都会低于概率得出的理论值。我惊讶人们这种挑选坏股票的天赋！对于交易了十年以上时间的我，仍旧无法认清其中的全部机制。股民都是抱着买入黑马的动机进行选股，而结果却是大面积地输给大盘。后来我曾经不止一次地建议新股民，不要选股，而直接买入ETF指数基金。我保证这能够让你跑赢绝大部分股民。而我在应对从股票池中选股的问题时，也不得不采用更大程度上的分散投资。对于交易而言，你包容掉的东西越多，你的战绩就会越好。当你舍弃了买入黑马的可能时，你就能得到大盘的平均升幅。如果你追求极致，那么就必须涉及更多的问题。而涉及的问题越多，你那丑陋的本能就越容易跑出来找麻烦。

很久很久以前，我就找到了一种能够获利的交易方法。但我并不甘心完全抛弃思考，而使用这种冰冷的规则。虽然我以这种方法作为全部交易的背景，但仍然总是想着超越！超越！再超越！回头看看这十年多的战绩，我“超越”没了数千万的利润！而这种结果，系统只不过是打击了我三次。

2001年的崩盘中，我因为忽略了一次卖出信号而亏损了将近50％；2003年我与市场赌气地追求确定性又亏损了将近50％；2006年我又在上证指数1800点时放弃了一次买入信号。就是这三次“超越”，让我丢掉了数十万（不考虑复利）。站在今天的角度来看，如果能够规避这三次自作聪明的操作，恐怕我的资产将会增加数千万。没错，就是数千万，这是十倍于我现在的财富。复利效应不但可以放大我们的利润，同样也能够放大我们的错误。这或许告诉我们，曾经所犯下很小的错误，在十年以后都会放大到你难以想象的地步。如果你的亏损来自于系统，那么除了接受而别无他法。但如果你的亏损来自于背离系统的错误，那么你恐怕在十年后会“付出”无法想象的代价！无论将来我多么富有，这三次错误都让我失去了十倍于我的财富。

回想我犯过的这三次错误，第一次我找到了“学习效应”来解释；第二次我下定决心放弃确定性；第三次更悲壮，我决定放弃微观上的思考。好在让我欣慰的是，最后一笔学费交得还算划算，毕竟我从期货市场中也获得了数倍的利润。但遗憾的是，我在进入期货市场以后，仍旧丢掉了一大笔豆油的利润。否则，我在期货市场上赚取的利润，要比我在股市上丢掉的利润大得多。更具有讽刺意味的是：除了这三次以外，我几乎都是完全按照信号操作。可能其间也成功地超越过几次，但带给我的利润（逃脱的亏损）也只不过是几万元罢了。

墨菲法则再一次应验。相信我，人类的思考在市场面前是负值的。如果人们绝对不思考，那么在随经济长期增长而不断上扬的股市中，绝对应该是获利人数多于亏损人数。这种由思考带来的“墨菲法则”的确很少引起人们的重视。《克罗谈投资策略》中说到，交易获利的关键在于对自身的约束。而书中讲述的一致性原则，事实上就是在暗示大家放弃微观上的思考。

我非常认同，交易系统能够获利的关键在于，必须放弃思考而坚持一致性原则。交易系统可以使能够获利的交易者获利更多，但却无法对亏损的交易者产生任何好处。请大家相信我：交易是一门朴实无华的艺术，而不是一门追求极限的艺术。交易讲究包容、低调、容忍、承担、放弃……本来交易是一种零和游戏，但人类的思考将其变成了负值游戏。想想看，放弃捕捉黑马你就能够与大盘同步。事实上，交易者全部的灾难，或多或少都来源于“不甘心”三个字！最后我终于并不情愿地相信：在庞大的市场面前，渺小的我真是一钱不值。

除了能够被我们找到、并且编制成为可量化系统的心理偏向以外，我们的任何本能反应仍旧是负值。简单地说，既然系统之外的本能反应是负值，那么我们就必须坚持一致性的规则。我并不反对大家思考，但这种思考应该是对系统负责的宏观思考，而不是对于微观上某笔特定交易的思考。实战中的任何疑问都需要暂时放下，而实战之后冷静下来再考虑是否应该修改系统的问题。当然，对于这个问题的思考一定不能“就事论事”地只考虑刚发生过的交易，而是要在大尺度下进行充分的测试。

\subsection{优化与组合}

{\kaishu
环境是变化的，而越能适应过去环境的方法，就越不可能适应变化以后的环境。

\hfill ——金融帝国
}

交易技术上的核心思路并非很多，那么为什么学习交易会如此困难呢？如果将全部的交易技术进行分类，那么核心思路也不过就那么几种。如果你只想学习其中的一种，那么短时间内你就可以毕业。交易之路之所以漫长，在我看来，完全是因为“组合”与“优化”的不归之路。所谓优化，就是为同一核心思路的交易技术找到最佳参数和公式。所谓组合，就是将不同核心思路的交易技术融合在一起。

对于趋势跟踪系统而言，核心思想是极其简单的。就是所谓的“亏损兑现、盈利挂起”，而做到这一点就需要考虑关于参数的问题。亏损达到什么程度我们应该兑现呢？或者说，趋势中的反向波动达到什么程度，我们可以视为反转呢？哈奇10％投资法则就能用最简单的方式回答这个问题：如果价格反向波动10％，我们就可以视为趋势的反转。或者说，用两条特定参数均线的交叉作为趋势跟踪尺度，同样可以回答这个问题。

但有一个不容忽视的问题，趋势跟踪系统必然会产生低正确率，甚至是“买天花板、卖地板价”的风险。几乎没有人心甘情愿地接受这种成本，那么“高智商”的人们就想通过“优化参数”或“优化公式”来绕过这种成本。自此趋向类指标的大家族中，就不断涌现新的成员。简单的问题在完美主义者的驱动下开始变得复杂起来。

我也不例外，同样花费了数年的时间，希望攻克避免假信号的终极课题。在应对某次假信号，使我止损在地板价时，我发现如果指标迟钝一点，那么就能够过滤掉这种假信号。我努力地寻找假信号的过滤器，来规避这种风险。所谓过滤器，就是在指标发出后，必须符合另外的一个条件才开始操作。事实上，这种过滤器的本质，只不过是将指标参数修改得迟钝一些罢了。或者说，某参数的指标加上某过滤器，同将该指标的参数变得迟钝一些是等效的。当然，如果过滤器的思路来自于趋势之外，那又会使跟踪信号变得不必然。而我所追求的是，只要有趋势出现，那么我一定要身在其中。从某种意义上来讲，无论是什么样的趋向指标，它们都只不过是对“一定程度”进行表述。甚至我们可以将不同的趋向指标，通过参数的选择来达到相同的“敏感当量”。

当我使指标迟钝一些后，就能规避当时已经出现的假信号。但迟钝后的指标，不久仍旧会出现止损在地板价的案例。这该如何是好呢？继续添加过滤器，使指标更迟钝，或者说更难发出信号吗？再过一段时间，我在一轮趋势行情中发现：过于迟钝的参数使得“掐头去尾”后利润极其有限；而如果采用敏感一些的参数，那么将会获利丰厚。

好在我比较务实，而没有在“头疼医头、脚痛医脚”的怪圈内“狗咬尾巴”地原地转圈。我认为根据一次特定的行情来改变自己的参数是愚蠢的，而必须建立一种长期观念来解决参数的问题。我意识到，继续玩这种“参数游戏”是绝对没有出路的，当我规避了一次假信号的问题时，却可能在历史走势图上又产生原本不会出现的新问题。我不能像“学习效应”中所说的那样，在敏感与迟钝的两极间摆来摆去，而应该给我的跟踪尺度找一个固定的“家”。

根据自己长期交易的经验，以及《混沌理论》《临界》等著作给我的启发，让我意识到“敏感与迟钝的等效性原则”。越敏感的跟踪尺度就能产生越多的假信号，但同时每次假信号的成本就越小，并且在趋势行情中就越能卖在距离最高点更近的价位；而越迟钝的跟踪尺度就能产生越少的假信号，但同时每次假信号的成本就越大，并且在趋势行情中就越会卖在距离最高点更远的价位。这种结论与《混沌理论》中的“本体相似性”不谋而合，同时也与《临界》中的“能量分布法则”别无二致。

至此，我们根本无法解决趋势跟踪系统的缺陷。那么参数的本质究竟是什么呢？按照《临界》的说法，复杂的现象并不可能具有常态分布的特征，而只会符合“能量分布法则”。也就是说，无论是趋势还是调整，都不可能存在典型的幅度。我们所能知道的不过是，幅度越大的趋势、整理所能够出现的次数就越少。由此看来，市场的本质是根本就没有参数的，而我们又必须用具有参数的技术来应对没有参数的市场。进入期货市场后，我产生了一个重大的转变，从对市场的探索，转变为对自己内心的探索。所谓参数的价值，并非来自于市场，而是来自于我们自身。趋势跟踪系统中的参数，根本就不是什么市场预测工具；而只不过是符合我们内心需求的资金管理工具。因为我们的环境不同、内心需求不同，所以参数的选择也是不同的，仅此而已！更忌讳亏损次数的交易者，可以选择更迟钝的尺度；更忌讳亏损幅度的交易者，可以选择更敏感的尺度；什么都忌讳的交易者，就可以离开市场了。由此看来，交易是一项很个性化的事业，选择的优劣是需要因人而异的。

在《通向金融王国的自由之路》中，作者介绍了很多衡量系统优劣的指标。我在看来，还有一个更为重要的指标，那就是系统的稳定性！我们可以从历史统计中，验证一种交易系统产生的战绩，但稳定性又该如何衡量呢？在我看来，衡量稳定性的方法有两种。第一，就是验证更久远的行情走势，是否也能产生同样的效果？第二，就是验证其他市场、其他品种的行情走势，是否也能产生同样的效果？

这就是我要说的交易系统的市场广度。在计算机非常发达的今天，我们可以轻松地找到在过去数年内最佳的获利方法。但这种获利方法是否在以后还能够有效呢？市场都是在不断变化的，那么变化后的市场是否同样能够保持这种完美呢？在我看来，能够在越多的市场产生效果的交易系统稳定性就越好；而只为特定市场、甚至特定品种而设计的交易系统就必然无法保持良好的稳定性。我们不难得出这样的结论：稳定性越好的系统就越依赖于理念；而稳定性越差的系统就越依赖于参数的优化。好的交易思想是应该在任何参数下都能有效的。如果一种交易思想只能在特定的参数下有效，那么它只可能是属于“微观规律”，并且会被市场快速淘汰。《通向金融王国的自由之路》说的一点也没错，越是过度优化的系统就越难在未来获利。由此我们可以得知，过度优化系统参数是绝对没有任何意义的。在我看来，对于一套交易系统而言，稳定性是比期望收益更为重要的一个指标。毕竟我们交易的是未来而不是过去，稳定性基本上就能代表系统在未来的有效性。对过去完美，而对未来无效的交易系统，我们要他干什么？

既然这种相同核心思维下的优化游戏不能解决缺陷，那么不同核心思维下的技术组合是否能够有好的效果呢？无论我们采用什么样的交易方法，交易方法也都不过是基于概率下的产物，而并不可能指向确定的结果。这是一个大家并不愿意接受的事实，但如果单一交易方法能够指向确定的结果，那么我们又何苦去组合交易方法呢？如果说，交易方法是基于概率下的产物，那么有效的时候就是收入、无效的时候就是成本。请注意，我们在将不同交易方法进行叠加的同时，不但叠加了收入，同时也叠加了成本。我们将众多不完美的交易方法组合在一起可能产生完美的交易方法吗？按照“圆木桶理论”“小蝴蝶的故事”“损耗效应”，我们可以得知组合交易方法的优势，并不能得到我们预期的效果。

在测试版的反馈中，我得知很多朋友仍旧坚信能够开发出双高系统。双高系统是指，正确率大于50％，与此同时赔率大于1。其实，我的本意从来就没有认为双高系统不存在，只不过是觉得对于双高系统的追求没有意义罢了。就我个人而言，即使不依赖于任何交易系统，单靠主观来交易就可以得到双高的战绩。一位想通过股市赚点钱的朋友向我寻求建议。我的答复是：要么你投资于期货市场，并且可以接受一定幅度的资金回撤；要么你就等着我告诉你股市中“千载难逢”的交易机会。不用说千载难逢，对于两三年一遇的交易机会来说，我都有很大的把握能够赚到不小的利润。我相信不单是我，很多读者也有这种能力，甚至不断亏损的交易者也具备这样的能力。问题是，很少有人能够耐得住性子，等待自己有把握的市场机会再进行交易。但为什么对于过滤交易系统的信号，更多的人却又乐此不疲呢？另一个问题是，看懂的利润我需要赚到，那么看不懂的利润我就不赚了吗？

我们可以在一套简单的系统上加入更多的过滤器，在《通往金融王国的自由之路》中也将其称为设置滤嘴。比如，加入250日长期均线作为过滤器，均线之下忽略买入信号，均线之上忽略卖出信号。或者，我们还可以找出其他更加有效的过滤方法。问题不在于过滤以后得到的战绩输出是否能够达到双高要求，而在于我们过滤出去的交易机会是否也能够使我们实现利润。哪怕这种利润并非十分可观，但我们又有什么必要放弃这些利润呢？毕竟，当你连续不间断地跟踪某合约的时候，即使你过滤掉了某个信号，闲置的资金也不可能挪作他用。因为也许明天就出现一个你需要执行的信号。

我们可以做一个不太恰当的比喻。对于企业而言，应该是追求最大的利润，而不是最大的利润率。那么只要是边际利润大于边际成本的话，企业就应该继续生产。对于一家手机店而言，可以销售手机、手机卡、充电器等等周边配套产品。当然，他们各自的利润率是不同的。甚至有的产品根本就不赚什么钱，而只是为了丰富产品的多样性，从而吸引尽量多的潜在客户。请注意，上述几类产品在销售过程中，并不存在任何冲突的可能性。难道店铺的老板应该为了追求更高的利润率，而将利润率低的产品下架吗？这或许让大家觉得可笑，但想象一下过滤信号的问题，大家还会觉得可笑吗？

如果一套交易系统一年可以产生10笔交易，同时4笔交易可以各自获得5元，而6笔交易各自会产生1元的亏损。你真的有必要通过自己的努力，将其中4笔亏损交易、以及1笔获利交易过滤出去吗？虽然这样你的正确率提高到60％，但你实际上却少赚了一元钱。

这还不是问题的全部，更多的过滤会让系统变得更加不稳定。如果你完全不进行过滤的话，就可以保证只要有趋势产生，你则必然身在其中。甚至还可以编制出一种趋势跟踪指数，并且研究这种指数的波动规律，从而运用于更高效的资金配置。虽然我没有办法证明，但我也能够感觉到市场中趋势特征与整理特征之间存在某种相关性。比如，整理的时间越长，就更容易产生大趋势。或者说，上下整理的次数越多，就越说明整理很可能即将完结。如果说趋势与整理存在某种相关性，那么这种相关性带来的稳定性甚至能够超越系统测试。比如说，趋势跟踪系统出现大亏损之前，往往都对应着非常快速的获利。这是一个明显的盈亏分布规律，但靠单纯的系统测试却无法发现这一点。很显然，虽然是相同的正确率与赔率，但趋势跟踪的盈亏分布，很可能与“抽球游戏”完全不同。

我相信，对于某单一核心思路的交易系统来说。边际利润永远大于边际成本，所以我们不应该放弃任何特定的信号。《海龟交易法则》中的统计似乎也在告诉我们，简单的双均线系统往往会有更好的表现，而更加复杂的系统并没有期望中的好效果。还是那句话，趋势跟踪形成的大赔率，在不断地嘲笑那些小亏损的同时，也在嘲笑我们规避亏损的努力。

如果希望开发一种双核心思路的交易系统的话，那么这实际上就等于是在使用两套不同的交易系统。两套交易系统在获利上，注定其中的一套没有交易价值，但或许获利小的那套系统存在分散价值。如果两套系统的成本付出期是不相关的、甚至是负相关的话，那么可以通过资金管理打造出更加平滑的系统。遗憾的是，我至今还没有找到趋势跟踪理念之外，其他可行的交易理念。我并不否定其他人进行这方面的努力，但还是建议将其视为两套简单的系统。如果实际证实其中一套系统不具备正期望收益，那么还是忘了这个主意吧。

测试版的另一个反馈是，很多朋友并不认同参数的敏感与迟钝具有等效原则。关于这个问题，我无法通过大量的测试来证明什么。但我曾经看到过一位在美国留学的网友推荐给我的美国股市测试图。这幅测试图非常形象，用X轴和Y轴代表两个均线的参数，而用颜色表示测试后得到的收益水平。我无法考证这幅测试图的出处，但表达的结果与我的感受基本相同。该图总共分为三组，并且每十年做为一组。我们可以发现，市场中统计出的最佳参数并非固定，而是不断地向右上角漂移。也就是说，随着时间的发展，趋势跟踪的最佳参数在逐渐趋于迟钝。如果我们对中国股市进行简单的测试，也不难发现这样的规律。这种现象或许说明，趋势跟踪的进入门槛将会越来越高，而收益率水平可能会不断下降。最低限度，这种统计似乎也证明了，并不存在某种特殊的最优数值。

质疑敏感与迟钝等效原则的另一个理由是：如果本体相似性是存在的话，那么适用于周线级别的趋势跟踪系统，同样也应该适应于日内级别。而事实上，我对日内的波动几乎完全无能为力。我曾经也因为这个理由而质疑过“本体相似性”，但后来我似乎找到了这个问题的一些线索。日内波动在趋势出现时，往往滑移差价会突然变得巨大。如果你的交易手数稍微大一些，成交就是一个明显的问题。巨大的滑移差价再加上手续费，很可能已经大于趋势跟踪系统所能够带来的利润。或许，杰西·利物默的日内交易在空中号子中能够获利，但在真实的市场中却只能亏损，正是因为这个道理。出现这种现象的唯一解释只可能是滑移差价，而只有顺势交易对这一点才更加敏感。我猜测，杰西·利物默也许仍旧使用曾经在空中号子的交易策略，只不过是把交易尺度转换为更大的级别。通常来说，让一个人彻底改变他的交易习惯是非常困难的。如果让我做整理交易，那么一定会表现得极其不自然。

\subsection{交易系统必定是个性化的}

{\kaishu
从某种意义上来讲，为自己选择一套适合的系统，就是一个不断进行利弊权衡的取舍过程。

\hfill ——金融帝国
}

《通向金融王国的自由之路》的开头就指出：设定目标是系统工作的主要部分。当时我并没有进入期货市场，自然对这句话就没有更深刻的理解。要知道股市里并不需要过多地考虑资金管理，所以使得这个“设定目标”的重要性就大大的弱化了。

对于目标不同的交易者而言，应该选择的系统也是不同的。即使是同样期望收益的交易系统，也是可能通过无数种不同的组合来实现的。从某种意义上来讲，为自己选择一套适合的系统，就是一个不断进行利弊权衡的取舍过程。如果你想通过期货市场，在一年的时间内解决自己的住房问题，那么与我希望通过期货市场实现复利效应而改变命运所使用的交易系统就会截然不同。即使这两种系统的期望收益是相同的，甚至正确率、赔率等都是相同的，但无论如何稳定性都会不同！不言而喻，快速获利的方法就必须承担更多的东西。还记得吗？我们从来都不曾真正得到什么，我们得到的只不过是通过付出另一些东西而换来的。

2006年我在期货市场上用两个月的时间获利了250％。如果单从结果上来看，这是一笔出色的交易。但如果站在交易系统的角度上来看，这种做法真是够得上自杀式的豪赌。毕竟，这种暴利是建立在近乎于满仓的条件下实现的！虽然很多可能毁灭我的因素并没有出现，但这并不等于没有出现的可能。如此快的获利速度，必然具有巨大的不稳定性。很显然，相当长一段时间之内，我都对那次疯狂的行为持否定的态度。甚至，发誓今后再也不使用这种自杀式的豪赌手法。但在多年后回顾那段经历时，我似乎已经越来越意识到那次豪赌的合理性。考虑到当时生活中的窘境，似乎只有背水一战才能解决实际问题。事实上，那时由于住房的问题，我家的财务状况几乎到了“资不抵债”（除了房产以外，账户中的资金的确已经小于负债）的境地。客观地说，我是一个追求公正并且拥有责任心的人。我不想平白无故地接受亲戚的资助，而把这种资助视为是自己的债务。几年以后我的确把这笔债务偿还了，但当时那笔债务对于我来说却是一个“天文数字”。当然，我可以选择“合适憨厚”地接受亲戚的资助。但这样一来，就永远找不回我的骄傲与自信。从另一个角度来讲，正是因为我实现了丰厚的利润，才让我有资格在今后的交易中实施更加稳健的交易策略。甚至还能够让我完全消除环境给我带来的压力与干扰。毕竟，在我屡战屡败的时候，没有哪个关心我的人不认为我是在不务正业。

我经常在网上遇到交易者向我寻求一些建议。当然，这种建议并非是做多还是做空的问题。对于抱着不劳而获心态的人，我不会深入地进行交流。我越来越发现，帮助一位陷入困境的交易者，找到一套适合自己并且可行的方案是我能力之外的事情。虽然我很理解他们所面临的困境，毕竟我也曾经历过相同的困境。对于交易的执迷或许已经使他们丢掉了现实中的一切，包括稳定的收入来源，而必须通过极少的资金靠交易养活自己。先不说太少的资金无法通过大面积分散实现获利的稳定性，即使交易能够实现获利，那么生活成本的支出也无法使其实现资金的积累。我只能无奈地承认，即使是我自己，也没有必然把握一定能够摆脱困境。难道，我应该劝他们豪赌一把吗？或许，我只能够建议他们找一份工作。毕竟，工资就能实现相对于账户中资金不小的收益率。此外，我并不赞成在金融领域找一份工作。道理很简单，没有人能够无耻到活在欺骗当中而没有道德压力，而只能是先骗了自己再去做自认为是帮助别人的事情。我的意思是说，从事金融行业必然会影响到你对市场本质的一些看法，而这种影响无疑对交易获利而言是绝对负面的。问题是，多年执迷于交易的结果，已经使得他们丧失了其他的谋生技能。

就我自身而言，我控制着足够实现分散投资的资金，并且这些资金完全是没有任何外在压力的自有资金。甚至当我偿还了父母的积蓄以后，没有任何人有权对我的交易指手画脚。要知道，在我上学的时候，即使用爸妈的资金也会让我感到压力巨大，甚至使交易完全走板！我相信，大家根本无法想象，爸妈看着我三次无效止损后，会给我带来多么大的压力。或许，正是因为自有资金的性质，能够让我接受更大的潜在资金回撤幅度。不难想象，由于资金的来源不同，很可能就需要采用不同的交易系统。另一个方面，我有一份与金融行业无关的稳定工作，并且继续过着简朴的生活。虽然我的工资收入相对于交易利润根本不值得一提，但我的确可以只靠工资来支付所有的生活成本，甚至还能够有些结余。这里顺便说一句，我之所以去上班，一方面是因为希望过相对正常的生活，另一方面就是希望摆脱对交易的依赖。如果我能够完全不依赖于交易，那么才更容易把交易做好。我始终相信，有资格拥有的人，一定是能够放下的人。

虽然我能够与困境中的交易者产生一种亲切的共鸣，但我真的完全无能为力！毕竟，没有人有资格代替别人作出选择，而真实的世界中也并不存在最优的正确答案。我最多也不过是帮他们分析各种选择的利弊。很显然，照搬我的交易系统对于很多人来说是根本不可能的。我曾经同一些管理他人资金的交易者聊天，建议他们应该对亏损有更大的包容程度。事实上，他们只要能够接受稍大一些的资金回撤幅度，那么战绩应该要比我好很多。但他们的角色就决定了，他们宁可放弃更多的潜在利润，也必须努力避免稍大一些的资金回撤。我突然发现，往往对错的问题，只有站在特定的前提下才有意义。最终，我仍旧赞同他们的努力方向是对的。不夸张地说，每一位交易者的风格，无不打上现实环境的烙印。

除了交易者的入市动机是不同的以外，每个交易者的心灵也是不同的。有的人把面子看得比生命还重要，那么他就不能使用过于敏感的趋势跟踪系统。甚至他都不能使用趋势系统，而不得不选用其他高难度的交易方式。说句实话，我现在使用的期货交易系统，就不是理论上效率最高的系统，甚至我的效率还达不到理论最佳值的一半。为什么呢？因为我的心灵要比电脑的中央处理器脆弱得多，我可能无法接受过大的资金回撤幅度。有时候我都不得不感叹自己是在股票市场上最激进的股民；而又是在期货市场上最保守的期民。在我看来，大多数的期货交易者，所持有的仓位都重得近乎于疯狂！理论上来讲，我可以把期货交易中的仓位提高一倍，但为了照顾我那颗“脆弱”的心灵而并没有那么做。通过电脑进行系统测试的价值是无可厚非的，但更有价值的却是通过长期实战而进行的心灵测试。我能够接受的最大潜在资金回撤幅度是40％（这只是一个极端值，并且采用了谨慎与稳健的原则），但我很难说对于其他人而言不会使其在心理上崩溃。要知道，资金最低点只有事后才能够知道。

事实上，交易者身上仍旧必然地存在一些“无效率”的风格。我们很难说，这种无效率的风格是没有任何现实意义的。我们可以举一个最简单的例子，我对于浮动盈利加仓的态度就非常谨慎。要知道，在破产风险可控的前提下，实现快速获利必定要依赖于对浮动盈利的使用。但对于我“永不丢掉”位置的根本原则来说，就不太适合使用浮动盈利加仓。道理很简单，使用浮盈加仓的交易者，绝不可能像我一样承担巨大的价格回撤，从而“证明”趋势的反转。我更多的是在追求“只要市场有趋势，我就必然有利润”的稳定性。这种得失的取舍，我们很难说什么才是更加划算的。或许，当我的系统在丧失一部分效率的同时，又在宏观上增加了更大的稳定性。

也许读者朋友们相对于我的思想而言，更想得到的是我的方法。我可以告诉大家，我使用的方法要比大家想象的简单得多。要知道这么简单的方法，在没有足够强大的理念作为后盾的前提下都是很难做到的。这里我提醒大家，不要把交易系统想象得过于复杂。从某种意义上来讲，“随便买、随时买、不要卖”的指数化投资也是一个不错的系统。如果真的要设计一套系统的话，不考虑自身环境以及需求是绝不可能的。甚至对于不同资金量的交易者都不可能“一视同仁”。这需要考虑市场容量的问题，以及可以进行风险分散的程度。很显然，对于索罗斯和巴菲特那样的重量级人物而言，完全不可能像我一样靠试错来选择市场方向。甚至，对于计算机水平超棒的交易者，也完全有可能开发出更加细腻的系统，而不用一定像我一样坚持简单。当我做交易的时间越长时，我就越懂得尊重他人的风格。市场上从来就不会有唯我独尊的方法，而越成熟的交易者就越相信这一点。

\subsection{如何开发一套适合自己的交易系统}

{\kaishu
问题是永远都解决不完的，从而探索的终点必将是我能够接受什么样的缺陷。

\hfill ——金融帝国
}

在测试版中，本章节的题目叫做《通往金融王国的自由之路》。对于希望系统学习系统开发的朋友来说，这的确是一本不容错过的好书。而对于趋势跟踪感兴趣的朋友，同样应该阅读一下《海龟交易法则》《趋势跟踪》这两本书。

我不得不承认，相对于交易感悟而言，系统地讲述如何设计交易系统的问题要困难得多。我相信市场上已经不存在任何新鲜的东西，如果我再不能从新鲜的视角来审视陈旧的问题，那么照抄《通往金融王国的自由之路》是否有任何意义呢？接下来，我还是尝试着通过自己的视角，来谈一谈关于交易系统的概念上的问题吧。

前面已经谈了很多交易系统的问题，那么交易系统又究竟是什么呢？在我看来，首先要消除大家对于系统的神秘感。从某种意义上来讲，每位交易者都拥有一套自己的交易系统。当然，这种交易系统未必是能够盈利的，甚至未必是可以量化的。但无论如何，我们每一位交易者的交易行为都绝不可能是杂乱无章的。从广义上来讲，交易系统不过是交易者在交易行为中所表现出的一些具有再现性的倾向。比方说，有的人就是喜欢买入市盈率低的股票；而有的人就是喜欢买入大幅下跌后的股票……无论这些倾向是否具有价值，每位交易者都具有某些习惯性的交易方式。无论交易者是否能够意识到自身的交易系统，这种交易习惯都会体现在真实的交易行为中。有意思的是，对于不同的交易者而言，交易习惯的确有某种相似之处。

对于相似之处来说，更多的体现在人们拥有相似的本能。我们暂且将形成于市场之外的系统称之为“本能系统”。问题在于，本能系统并不能使交易者实现利润。由于本能系统的无效率，就使得我们有必要开发出具有正期望收益的交易系统。很显然，后天开发的系统，就是一种对本能系统的约束。交易系统对于交易行为的改变，已经很类似于生活中通过纪律对坏习惯的改变。

我们可以从两个角度来开发系统。一种方法是对市场波动特征进行大量的统计，从而找到市场波动中的偏向性。当然，这需要一些我并不具备的计算机知识。另一种方法就是假设市场是人性的一面镜子，从而用对本能偏向的研究来代替对市场偏向的研究。至少对于我来说，后者显得更加容易一些。但问题是，这需要更多的交易经验，以及对自己内心冰冷的审视。

当我们找到了市场的偏向以后，就需要用可量化的方式将其表述出来。对于不同的交易者而言，所使用系统的自由度很可能是不同的。比方说，在期货市场中更容易使用刚性的系统，而在股票市场中就需要通过一些主观的方法来解决选股等问题。我们很难说，越刚性的系统就越好，或者越不好，但我仍旧认为初学者应该使用需要更少主观判断的系统。交易系统的本质并不深奥、也不神秘。从某种意义上来讲，它很类似于减肥或者戒酒。当我们意识到生活中某些不好的习惯需要改正时，就必须靠制定一些纪律来约束坏习惯。很显然，越简单的纪律其可操作性就越强。否则，我们就会不自觉地将纪律变通，从而多吃饭，或者多喝酒。对于交易也是同样的道理，越具有弹性的系统，就越难约束好自己的本能。

交易系统可以很简单，也可以很复杂，但并非是越复杂就越好。对于我的系统而言，就是一套简单的双均线系统，但使用起来效果却很不错。说实话，本人的交易系统并非是有意开发出来的，而是通过长期的实战而磨合出来的。当我在交易中遇到某个不能接受的问题时，就会思考解决之道。请注意，为了解决一个老问题的同时，必然会增加一些新问题。总之，问题是永远都解决不完的，从而探索的终点必将是我能够接受什么样的缺陷。有的时候，转了一个大圈子后发现，我需要解决的问题却是我必须接受的问题！而其他的方法，只不过把这种缺陷转换为另外一种形式的缺陷罢了，甚至后者的缺陷更加难以接受。具有讽刺意义的是，我大部分努力的结果，就是发现更多的等效性原则，从而使交易系统不得不又回到原点。

我的交易系统需要的主观判断要比大家想象的更少。但无论如何，交易都很难实现“无人值守”。即使是期货市场，微观上的资金管理、市场选择都需要靠主观判断来实现。对于资金管理而言，我最多只能规定不能如何，而不可能时时回答出目前的仓位应该是多少。或许，交易战绩比拼的正是那些不可量化的东西，但初学者应最大程度地回避不可量化的东西。要知道，系统只不过是我们改变坏习惯的工具。而当我们的坏习惯已经完全改正以后，系统就是免去它存在的意义。我经常戏称，交易者成长的道路，很类似于达芬奇画鸡蛋的学习过程。随心所欲不逾矩，只是我们追求的终极目标，而在此之前只能够靠更多的自我约束来走完成熟之路。

我最多只能够谈一谈原则性的问题，没有人能够代替你进行交易中的取舍。\textbf{系统的开发之路，就是一条不断取舍之路，而只有你自己才能够决定应该追求一些什么，同时又能承担一些什么。}前面已经说过，我对系统的寻找过程，是一种发现问题、并且解决问题的过程。解决问题的前提是发现问题，而我并不知道你有什么问题需要解决。如果你的问题是如何交易才能永远避免亏损，那么我只能说这个问题是无解的，或者说只有离开市场才能永远不会出现亏损；如果你的问题是如何才能将止损的次数降到最低，那么我只能说设置更加迟钝的止损点；如果你的问题是如何才能够让自己的卖出价位距离最高点更近一些，那么我只能说设置更加敏感的跟踪止损。发现没有？解决一个老问题，的确会增加另一些新问题。考虑到兼容性的情况，你一个问题与我的答案不同，就有可能使其他问题的答案与我有天壤之别。有人说，市场上没有正确的答案，而交易者的差别就在于如何提出正确的问题。从一位交易者提出的问题，往往就能看出他在市场中的境界。既然我无法得知你的问题，同时又不能代替你决定该接受什么样的缺陷，那么我自然无法帮助你设计交易系统。当然，你的系统未必像我一样来自于实践的磨合，而可以是更多地依赖于软件进行的统计，或者是对心理学大量的研究。但无论你的系统是如何形成的，仍旧无法摆脱同自身的磨合过程，而这个过程恐怕要比大家的想象花费更长的时间。

悲观地说，没有人能够教会你操作一套实现盈利的系统，甚至没有人能够帮助你开发适合自己的系统。我不得不承认，《通往金融王国的自由之路》是对我影响最大的读物。我得知本书的存在也是非常具有戏剧性的。曾经在论坛上看到某位网友质疑该书中“正确率与赔率是互损”的观点，我突然发觉这种观点与我多年来的经验不谋而合。我决定详细研究一下该书，阅读本书的过程就是不断产生共鸣的过程。真的有一种相见恨晚的感觉，并且对我至今沿用的思维方式产生了重大的启发。我突然发觉：谁也不可能真正教会你什么，最多也只不过是唤醒你已经经历过，并且感悟到的内容。

除非你真正的经历过，否则任何书籍也不会给你带来帮助。我熟读数遍该书的内容，但还是在刚进入期货市场时满仓交易。也许实践真的是成熟起来的唯一方法！不过阅读优秀的书籍能够加速我们的成熟也是一个不争的事实。我并不奢望我的作品能够教会你什么，只要你能够在阅读中产生一些共鸣，并且得到哪怕是一点点的启发，那我已经心满意足了。

\subsection{昆虫系统}

{\kaishu
对于趋势跟踪者而言，我们不怕上涨也不怕下跌，甚至我们连整理都不怕。毕竟小的整理能够被我们所包容，大的整理能够让我们实现利润。请注意，我们怕的就是那种“恰到好处”的整理！

\hfill ——金融帝国
}

昆虫系统的名字来源于“恐龙效应”中讲述的道理，由此看来它是一个追求稳定性的系统。对于设计交易系统而言，期货市场需要考虑更多的问题。而更加灵活的交易制度，反而更容易设计出具有稳定性的系统。这里就和大家谈谈这种为期货市场设计的“昆虫系统”吧。

昆虫系统的核心思想就是趋势跟踪系统，这绝对是完全平淡无奇的东西。我们永远在期货合约中保持自己的位置，简单地说就是：一定程度的上涨就平空翻多；一定程度的下跌就平多翻空；而无论是继续上涨还是下跌都必须无条件持仓。所谓“一定程度”的确定就需要找到一个具体的参数。在我看来，只要能够包容隔夜跳空风险的跟踪尺度都是可行的。这里需要指出的是，跟踪方法需要直观，而不要自己给自己找麻烦。对于敏感程度而言，就完全是一个因人而异的问题了。

风险控制采用大面积分散策略。这样一来我就完全依赖概率优势，而不是任何特定的行情。我相信，如果我（的方法）是对的，那么十个我也是对的。这种大面积分散策略充分体现了“无为而治”的哲学思想。由于抛弃了思考，以及不依赖任何特定行情，所以该系统的稳定性绝佳。毕竟，能够使资金出现大幅度波动的偶然事件，并不太可能同时发生。理论上来讲，分散得越广效果就越好。我们可以将资金分配在农产品市场、工业品市场、股指期货市场，并且在各自市场的品种内再进行更细致的分散。

资金管理采用等额投资法。单从资金管理上来思考的话，那么赢冲输缩的下注策略具有明显的优势；但从趋势跟踪系统的特点上来思考的话，趋势行情通常产生于整理行情过后，也就是说连续几次亏损后才能出现大利润。这样一来输冲赢缩的方式更具有优势。我不希望看到最后产生利润的开仓数量是最小的；更不希望陷入亏损的泥潭而无法自拔，甚至最终被市场扫地出门。那么最好的办法就是等额投资法。事实上，本系统在设计上已经体现了赢冲输缩的思维方式。

具体的操作方式就那么简单，下面我说一下这种系统的原理。对于趋势跟踪者而言，我们不怕上涨也不怕下跌，甚至我们连整理都不怕。毕竟小的整理能够被我们包容，大的整理能够让我们实现利润。请注意，我们怕的就是那种“恰到好处”的整理！

对于相同品种的不同月份合约而言，它们必然会存在走势上的差异，这是不容置疑的！这时我们甚至可以在品种的不同月份上进行建仓。道理很简单，因为不同月份间的走势具有差异化，那么整理幅度就绝对不可能同时都“恰到好处”。可能一个月份发出了假信号，而另一个月份没有信号产生；也可能一个月份买在了天花板上，而另一个月份虽然也出现了假信号，但只不过让我们亏损了一点罢了。这样一来，最低限度我们不可能同时出现“天花板买入与地板价卖出”这样极端“荒唐”的交易。对于不同品种而言，走势的差异性将会更大，那么自然就会得到更好的效果。

对于整体农产品市场而言，它们具有一定程度的相关性。对于农产品强大的趋势运动而言，它们经常表现为齐涨共跌。但对于整体农产品的整理行情而言，它们很可能出现分化。如果我们用完全不思考的方式，将资金分散在豆系、玉麦、糖棉这些农产品合约上会出现什么现象呢？

如果市场都不知道未来会如何的时候，我们又何苦非要看穿市场的将来呢？这种方法妙就妙在：市场有序我就有序，市场无序我就无序。事实上，这就是传说中的以天合天。当农产品市场整体出现整理行情的时候，我们可能是在不同品种间、甚至是不同月份间进行多空的无序持仓。这时的风险对冲，基本上就可以让我不惧怕国外市场夜间大幅度的波动。而如果农产品整体走出趋势行情时，我的持仓方向就会陆续调整一致。只要是趋势继续进行，那么我就可以获利丰厚。当然，这时由于持仓方向的相同，就可能出现大幅度的资金回撤。这里还有一个妙笔，那就是赢冲输缩的奥秘。当可能出现大幅度亏损的时刻，必定是我们已经有利润的时刻，或者说是我们最亏得起的时刻。而在整理行情中，我的多空持仓是无序的，这时就不可能出现大的资金回撤。更有意思的是，即使整体农产品市场是整理阶段，但个别品种也会表现出不同的趋势特征。比如2007年上半年，豆油的多头就可以获利丰厚，而玉米的空头也可以获利丰厚。通过理论可以计算得知，当一个品种走出不那么壮观的趋势时，就能立刻对冲掉两到三个其他品种上高买低卖左右挨打的成本。

把系统进行历史统计，然后计算得出各种参数。我们不难发现：正确率为40％，赔率为3倍以上；资金可能出现的最大回撤幅度为30％；18个月后出现亏损的概率小于5％；除极端趋势行情出现的年份外，年化收益率在30％左右，而出现极端趋势的年份收益率会超过100％，甚至更多；操作频率为平均每周一次，而全部持仓周转一次需要较长的时间，所以交易成本不大。对于单个品种来说，平均每年只会产生2.5对的交易信号。最为重要的是，该系统没有进行任何优化，甚至不需依赖任何特定的行情，所以该系统的稳定性极佳。

\subsection{交易系统的前景}

{\kaishu
从某种意义上来说，交易系统的本质就在于相信技术指标。

\hfill ——金融帝国
}

交易系统是最近一段时间才脱颖而出的新概念。从某种意义上来讲，它与上个世纪出现的技术指标并没有什么本质的不同。它们同属于统计学的范畴，甚至部分技术指标被循环使用就能组成一个简单的交易系统。当然对于那些只能给出提示作用，但并不能循环使用的技术指标而言，交易系统还是与其有不同之处的。由此看来，交易系统把用于分析的技术指标挡在了门外，而给予交易者循环使用技术指标的一些信念。

从某种意义上来说，交易系统的本质就在于相信技术指标。这种包容缺陷、宏观评价、尊重技术的宝贵思想确实具有跨时代的意义。在我看来，交易系统的精髓不在于“系统”，而在于“交易”。也就是说，交易行为能否产生良好的效果，并不取决于系统的本身，而是取决于交易者的自身状态。我可以诚实地告诉大家，本人的系统就是单一的、常用的技术指标，但这与实现稳定获利并没有丝毫的矛盾之处。

对于作为新鲜事物的交易系统而言，它必然会经历其他新鲜事物所遵循的发展轨迹。就像上个世纪出现的技术指标一样，经历了陌生、接受、狂热、平稳、淡忘的历程。我曾经听到过这样的说法：“一位交易者谈论技术指标的频率越高，就说明这位交易者的水平越差。”事实上，大众对于技术指标的关注程度已经大大小于十年前。是否交易系统在几十年以后也会出现同样的结局呢？

对于优秀的交易者而言，技术指标同交易系统是绝对等效的。武器是什么并不重要，重要的是使用者的能力。对于交易而言，无论是技术指标还是交易系统。起作用的并非是它们本身，而是使用者的信念体系、心境、动机、目标等。我的意思是，如果不从改变思维方式、约束本能入手，那么再好的工具都不能产生任何效果。这就是本章并不侧重于交易系统的开发、介绍等内容，而是更侧重于使用交易系统时需要注意哪些问题的原因。如果你不能认同并且拥有使用交易系统需要具备的要求，那么给你再好的交易系统都是没有任何意义的。但如果你能够认同并且拥有使用交易系统需要具备的要求，那么具体使用什么样的交易系统对你来说已经并不重要了。别人的东西再好也是别人的，请读者朋友们相信自己，从而找到最适合自己的交易系统。

\clearpage
